\documentclass[11pt,a4paper]{book}

\usepackage{fontspec}
\usepackage[ngerman]{babel}
\usepackage[a4paper,top=2.5cm,bottom=2.8cm,inner=2.6cm,outer=2.2cm]{geometry}
\usepackage{amsmath,amssymb,amsthm}
\usepackage{booktabs}
\usepackage{microtype}
\usepackage{parskip} % kein Absatzeinzug, stattdessen kleiner Abstand
\usepackage{emptypage} % Leerseiten wirklich leer (keine Kolumnentitel/Seitenzahlen)
\setlength{\emergencystretch}{1.5em} % Notdehnung nur für sonst überlaufende Absätze
\usepackage{horimetrik}
\usepackage[hidelinks]{hyperref}
\hypersetup{pdftitle={Horimetrik — Zahlen mit Horizont},
  pdfauthor={Christian Felix Bürckert},
  pdfsubject={Wert, Struktur und Horizont in einem Zahlensystem mit fallenden Basen},
  pdfkeywords={Horimetrik, Fakultätszahlensystem, Mischbasen, Halbring, Seitentreue}}

\title{Horimetrik\\\large Zahlen mit Horizont}
\author{Christian Felix Bürckert, M.\,Sc.\\[4pt]
\normalsize unter mathematischer Mitarbeit von Claude (Anthropic)\\
\normalsize und Ideenbeiträgen von GPT (OpenAI)}
\date{Arbeitsstand \today\\[4pt]
\normalsize Lizenz: CC\,BY\,4.0 (Text) --- Begleitprogramme: MIT}

\begin{document}

\maketitle
\tableofcontents

\chapter*{Vorwort: Was passiert, wenn man die Exponenten umdreht?}
\addcontentsline{toc}{chapter}{Vorwort}
\markboth{VORWORT}{VORWORT}

Dieses Buch ist nicht am Schreibtisch entstanden, sondern im
Kinderzimmer, beim Zubettbringen meiner Tochter. Ich lag
neben ihr im Dunkeln und wollte nicht einschlafen, nur begleiten --
und nichts hält den Geist so zuverlässig wach wie ein Gedanke, der
ihn nicht loslässt. Bei mir sind das Zahlen; mein Kopf rechnet von
selbst. An jenem Abend übte ich Kopfrechnen, aus dem etwas
trotzigen Vorsatz, das Rechnen im Zeitalter der Computer nicht zu
verlernen: 728 geteilt durch 4. Die 700 sind 7 mal 100, also 7 mal
4 mal 25 -- durch 4 geteilt bleiben 7 mal 25, das ist 175; die 28
liefern trivial 7; zusammen 182.

Nichts daran ist bemerkenswert, außer dem, was mir beim Zerlegen
auffiel: wie selbstverständlich ich Hunderter, Zehner und Einer
auseinandersortiere. Von dort war es ein kurzer Weg zu dem, was ich
über Stellenwertsysteme wusste -- Binärzahlen, das Zwölfersystem,
das man an den Fingergliedern abzählen kann --, und zu einer Frage,
die ich so noch nirgends gestellt gesehen hatte: Was passiert, wenn
die Zählbasen \emph{innerhalb} einer Zahl wechseln? Dass es solche
Systeme längst gibt -- die Fakultätszahlen --, wusste ich in dieser
Nacht noch nicht; das ergab erst die erste Recherche. Aber
aufsteigende Basen kamen mir langweilig vor. Ich wollte es
andersherum: ganz rechts die meisten
Möglichkeiten.\footnote{Genau genommen wird kein Exponent
umgedreht, sondern die Richtung, in der die Basen wachsen -- die
Gewichtsordnung des Stellenwertsystems. Der Titel folgt dem
ursprünglichen Gedanken, nicht der Terminologie.} Wer diesen
Gedanken ernst nimmt, stößt sofort auf eine Zwangsläufigkeit: So
kann man nicht bis unendlich zählen -- Unendlichkeit wäre hier
Unsinn gewesen, also: endlich. Was zunächst wie ein Defekt
aussieht, ist die Geburt des zentralen Begriffs dieses Buchs: Jede
Zahl lebt in einem \emph{Horizont}, einem endlichen Raum mit einer
festen Anzahl von Plätzen. Meine Tochter schlief da längst. Meine Frau
erzählte mir später, ich sei gut eine Stunde oben gewesen; gefühlt
waren es fünfundzwanzig Minuten.

Dass führende Nullen in so einem System die Zahlen verändern -- die
Merkwürdigkeit, die alles Weitere antreibt --, wusste ich in jener
Nacht noch nicht. Sie kam ans Licht, wie man Ideen heute prüft: Ich
setzte mich an den Computer und erklärte den Gedanken GPT, mit der
Bitte, nachzusehen, ob es das alles schon gibt. Das erste Urteil
war trocken: Unsinn. Aber wir blieben dran, und nach einiger Zeit
merkten wir gemeinsam, dass es kein kompletter Unsinn war. In
diesen Gesprächen tauchten die Nullen auf, die vorne an einer Zahl
plötzlich etwas bedeuten, und viele der Gedanken, die dieses Buch
aufgreift. Ein Satz aus dieser Zeit ist mir geblieben: 4 mal 5 kann
eine Fläche sein -- und wenn die 4 ein Skalar ist, eine Länge; in
beiden Fällen 20. Mit der Horimetrik hat das nur am Rande zu tun,
aber es machte Eindruck -- und findet in
Kapitel~\ref{ch:strukturpolynome} ein spätes Echo. Der Name
übrigens entstand ganz unfeierlich: „Horizont-Arithmetik“ war mir
zu lang. Also Horimetrik.

Irgendwann war das Thema zu groß für Plaudereien, und GPT und ich
waren uns einig, dass es einen eigenen Agenten verdient. So kam
Fable ins Rennen -- das Claude-Modell von Anthropic, das auf der
Titelseite als mathematischer Mitarbeiter firmiert. Ich sage es
ehrlich: Mathematisch wurde ich schnell abgehängt. Mein
Informatikstudium hat mir Mathematik beigebracht, samt Nebenfach --
aber das liegt über ein Jahrzehnt zurück. Die Grundzüge verstehe
ich, und die waren spannend genug: nicht, weil es eine Anwendung
gäbe, sondern weil das alles neuartig genug wirkte, um ihm
nachzugehen. Viele mathematische Felder waren zu Beginn ohne
Anwendung. Fable kann Mathe; also ließ ich alles nachrechnen und
eine saubere Theorie aufstellen, und mit der Zeit fanden sich immer
mehr Übergänge zur etablierten Mathematik. GPT steuerte
zwischendurch weiter Ideen bei. Beiden gilt mein Dank -- lesen Sie
selbst, was die KI und ich gefunden haben.

Das Buch verfolgt zwei Ziele. Erstens will es die Mathematik dieses
Raums ehrlich entwickeln -- mit Definitionen, Beweisen und auch mit
den Irrwegen, die wir unterwegs gegangen sind. Zweitens will es
lesbar bleiben: Jede technische Passage kehrt an die Oberfläche
zurück, bevor die nächste beginnt. Wer nur die formelfreien
Abschnitte liest, soll trotzdem verstehen, worum es geht.

Ob die Horimetrik eine eigenständige Theorie wird oder eine
besonders lehrreiche Spielerei bleibt, wussten wir beim Schreiben
selbst noch nicht. Genau davon handelt dieses Buch.

\chapter{Die Idee: Zahlen, die wissen, wie viel Platz sie haben}
\label{ch:idee}

% Ebene: Oberfläche. Keine abgesetzten Formeln, nur Zahlen im Fließtext.
% Entstehungsreihenfolge beachten: umgedrehte Basen -> Endlichkeit ->
% führende Nullen (Korrektur vom 13.07.2026).

\section{Eine Ordnung, die niemand hinterfragt}

Jedes Stellenwertsystem beruht auf einer stillen Vereinbarung. Im
Dezimalsystem zählt die hinterste Ziffer Einer, die davor Zehner, dann
Hunderter -- jede Stelle wiegt das Zehnfache ihrer Nachbarin. Im
Binärsystem dasselbe mit Zweierschritten. Und es gibt ein weniger
bekanntes, aber altehrwürdiges System, in dem die Basen nicht konstant
sind, sondern wachsen: das Fakultätssystem. Dort darf die hinterste
Ziffer nur 0 oder 1 sein, die nächste 0 bis 2, die übernächste 0 bis 3,
und so weiter. Je weiter links eine Stelle steht, desto mehr darf sie
zählen \emph{und} desto mehr ist sie wert.

In allen diesen Systemen zeigt dieselbe Richtung nach oben: Die Stelle
mit dem meisten Spielraum trägt das meiste Gewicht. Das wirkt so
selbstverständlich, dass man es nie ausspricht.

Die Horimetrik beginnt mit der Frage, was passiert, wenn man genau das
umdreht. Was, wenn ausgerechnet die \emph{engste} Stelle -- die, die
nur 0 oder 1 kennt -- das größte Gewicht bekommt? Wenn die Stellen
nach rechts hin immer großzügiger werden, ihr Einfluss aber immer
kleiner?

\section{Die erzwungene Endlichkeit}

Die erste Konsequenz kommt sofort, und sie ist keine Designentscheidung,
sondern ein Zwang. Ein gewöhnliches Zahlensystem wächst nach links ins
Unendliche: Wird eine Zahl zu groß, stellt man eine weitere Stelle
voran, mit noch größerem Gewicht -- Platz ist immer. Im umgedrehten
System geht das nicht. Links steht bereits die engste Stelle, die mit
der Basis zwei, und unterhalb von zwei gibt es nichts mehr. Das System
ist nach links \emph{versiegelt}.

Wer mit einer festen Stellenzahl arbeitet, bekommt deshalb einen
endlichen Vorrat an Zahlen. Mit einer Stelle sind es zwei, mit zweien
sechs, mit dreien vierundzwanzig, mit vieren hundertzwanzig -- jede
zusätzliche Stelle multipliziert den Vorrat, und die Größen, die dabei
entstehen, sind genau die Fakultäten. Ist der Vorrat aufgebraucht,
hilft kein Übertrag: Man braucht ein \emph{neues System} mit mehr
Stellen, in dem jede alte Ziffer plötzlich ein anderes Gewicht trägt.

Diesen endlichen Lebensraum einer Zahl nennen wir ihren
\emph{Horizont}. Was zunächst wie ein Defekt des umgedrehten Systems
aussieht -- man kann nicht einfach weiterzählen --, ist in Wahrheit
sein interessantester Zug: Jede Zahl trägt hier nicht nur einen Wert,
sondern auch die Information, in welchem Raum dieser Wert lebt. Die
Gesamtheit aller Horizonte bleibt dabei unendlich; nur jeder einzelne
ist endlich. Unendlichkeit gibt es weiterhin -- aber als Treppe, nicht
als Gerade.

\section{Der zweite Fund: die Nullen}

Die führenden Nullen kamen erst später, und sie kamen als Überraschung.
In jedem gewohnten System ist eine vorangestellte Null reine Kosmetik:
012 ist 12. Im umgedrehten System stellt man eine Null voran -- und
verändert damit die Zahl.

Ein Beispiel mit kleinen Zahlen. Im Horizont mit zwei Stellen bedeutet
die Ziffernfolge 12 den Wert fünf. Stellt man eine Null davor, steht
dort 012 -- aber jetzt in einem Horizont mit drei Stellen, in dem alle
Gewichte anders verteilt sind. Der Wert ist sechs. Noch eine Null:
0012, Wert sieben. Die Ziffern 1 und 2 haben sich nicht bewegt, und
doch zählt die Zahl bei jeder Verlängerung eins weiter.

Das ist der Moment, in dem aus einer Kuriosität ein
Untersuchungsgegenstand wird. Denn wenn eine führende Null den Wert
verändert, dann gibt es plötzlich \emph{zwei verschiedene Antworten}
auf die harmloseste aller Fragen: Was heißt es, einer Zahl mehr Platz
zu geben?

\begin{itemize}
  \item Man kann ihren \emph{Wert} festhalten -- dann müssen sich beim
        Umzug in den größeren Horizont ihre Ziffern umorganisieren.
  \item Man kann ihre \emph{Gestalt} festhalten -- die Ziffern bleiben
        stehen, eine Null kommt davor -- dann verschiebt sich ihr Wert.
\end{itemize}

Im Dezimalsystem sind das dieselbe Operation. Hier sind es zwei, und
sie führen in verschiedene Welten. Fast alles, was dieses Buch
untersucht, lebt in der Lücke zwischen diesen beiden Begriffen: die
beiden zunächst sichtbaren Identitäten einer Hori-Zahl (eine dritte
wird sich später zeigen), die Arten aufzuräumen, die Rechenwelten
und die Spannungen zwischen ihnen.

\section{Was hier eigentlich erfunden wurde -- und was nicht}

Ehrlichkeit gehört von Anfang an dazu: Die Bausteine dieses Systems
sind sämtlich bekannte Mathematik. Zahlensysteme mit gemischten Basen
gibt es seit Jahrhunderten, das Fakultätssystem ist gut untersucht,
und dass seine Räume Fakultätsgröße haben, weiß die Kombinatorik
lange. Neu -- jedenfalls haben wir es in dieser Form nirgends gefunden
-- ist die Kombination: die umgedrehte Leserichtung in einem endlichen
Raum, der Horizont als Bestandteil der Zahl, die bedeutungstragenden
Nullen und die beiden konkurrierenden Arten zu wachsen. Ob diese
Kombination eine eigenständige Theorie trägt, war die Wette, mit der
dieses Projekt begann. Die folgenden Kapitel legen die Karten offen.

\begin{oberflaeche}
Falls von diesem Kapitel nur ein Satz bleiben soll, dann dieser: Die
Horimetrik entstand nicht aus dem Wunsch, Zahlen anders zu schreiben,
sondern aus der Frage, was eine Ordnung erzwingt, wenn man sie
umdreht. Die Antwort war: Endlichkeit -- und aus der Endlichkeit
folgte alles Weitere von selbst.
\end{oberflaeche}

\chapter{Endliche Horizonte und Fakultätsschalen}
\label{ch:horizonte}

% Register: D1, D2, S1. Erst Beispiele, dann Definition.

\section{Ein Nachmittag in einer Welt mit 24 Zahlen}

Bevor wir definieren, rechnen wir. Wir nehmen drei Stellen. Die erste
darf 0 oder 1 sein, die zweite 0, 1 oder 2, die dritte 0 bis 3. Die
Gewichte laufen von rechts: Die letzte Stelle zählt Einer, die
mittlere zählt Vierer (denn die letzte Stelle wird mit ihren vier
möglichen Werten „voll“, bevor die mittlere weiterschaltet), die
erste zählt Zwölfer. Wer alle erlaubten Kombinationen der Reihe nach
aufschreibt, zählt von 0 bis 23:

\begin{center}
\begin{tabular}{cccc}
$000=0$ & $001=1$ & $002=2$ & $003=3$\\
$010=4$ & $011=5$ & $012=6$ & $013=7$\\
$020=8$ & $021=9$ & $022=10$ & $023=11$\\
$100=12$ & $101=13$ & $102=14$ & $103=15$\\
$110=16$ & $111=17$ & $112=18$ & $113=19$\\
$120=20$ & $121=21$ & $122=22$ & $123=23$\\
\end{tabular}
\end{center}

Keine Lücke, keine Doppelung: genau $2\cdot3\cdot4=24$ Zahlen. Bei
$123$ ist Schluss -- jede Stelle steht auf ihrem Maximum, und links
anbauen geht nicht, weil es unter der Basis zwei nichts gibt. Diese
Welt mit 24 Plätzen nennen wir den Horizont $H_3$. Wer die 24
erreichen will, braucht eine neue Welt: $H_4$ mit
$2\cdot3\cdot4\cdot5=120$ Plätzen, in der alle Gewichte andere sind.

\section{Die Definitionen}

\begin{definition}[Horizont; Register D1]
Für $n\in\Nn$ sei
$\Hn{n}=\{(a_1,\ldots,a_n)\mid 0\le a_i\le i\}$.
Die Stelle $i$ besitzt die Basis $i+1$; damit ist
$|\Hn{n}|=2\cdot3\cdots(n{+}1)=(n+1)!$.
\end{definition}

\begin{definition}[Wert; Register D2]
$V_n(a)=\sum_{i=1}^{n}a_i\,\dfrac{(n+1)!}{(i+1)!}$, gleichwertig
rekursiv $v_0=0$ und $v_i=(i{+}1)v_{i-1}+a_i$ von links nach
rechts, mit $V_n(a)=v_n$.
\end{definition}

\begin{satz}[Register S1]
$V_n\colon\Hn{n}\to\{0,\ldots,(n+1)!-1\}$ ist bijektiv.
\end{satz}

\begin{proof}
Induktion über $n$. Für $n=0$ ist nichts zu zeigen. Die Rekursion
liefert
\[
V_n(a)=(n{+}1)\,V_{n-1}(a_1,\ldots,a_{n-1})+a_n,
\qquad 0\le a_n\le n;
\]
Division mit Rest durch $n+1$ bestimmt also $a_n$
und den Rest der Ziffern eindeutig (Injektivität). Der größte Wert
ist $(n{+}1)(n!-1)+n=(n{+}1)!-1$, und da Definitionsbereich und
Zielbereich beide $(n+1)!$ Elemente haben, folgt die Surjektivität.
\end{proof}

Die Umkehrabbildung $E_n$ -- einen Wert in seine Ziffern zerlegen --
ist die fortgesetzte Division mit Rest durch $n{+}1, n, \ldots, 2$,
also genau der Algorithmus, mit dem man Sekunden in Tage, Stunden,
Minuten zerlegt, nur mit regelmäßig wachsenden statt willkürlich
gemischten Basen.

\section{Fakultätsschalen}

Jede natürliche Zahl $x$ hat ein kleinstes Zuhause: den
\emph{Minimalhorizont} $\mu(x)=\min\{n:x<(n+1)!\}$. Die natürlichen
Zahlen zerfallen dadurch in \emph{Schalen} $S_n=\{x:\mu(x)=n\}$ --
und die sind begrifflich etwas anderes als die Horizonte: Der
Horizont $\Hn{n}$ ist der ganze Raum und beherbergt die Schalen
$S_0$ bis $S_n$; die Schale $S_n$ ist sein \emph{Neuland}, also die
Werte, die erst mit ihm möglich werden.

\begin{center}
\begin{tabular}{lll}
Schale & Bereich & neue Werte\\
\hline
$S_0$ & nur die $0$ & $1$\\
$S_1$ & nur die $1$ & $1$\\
$S_2$ & $2$ bis $5$ & $4$\\
$S_3$ & $6$ bis $23$ & $18$\\
$S_4$ & $24$ bis $119$ & $96$\\
$S_5$ & $120$ bis $719$ & $600$\\
$S_6$ & $720$ bis $5039$ & $4320$\\
\end{tabular}
\end{center}

Die Schalengrenzen sind zahlentheoretisch markante Orte: Für
$n\ge1$ ist die untere Grenze der Schale $S_n$ die Fakultät $n!$ --
durch \emph{jede} Zahl von $1$ bis $n$ teilbar. Die Zahl direkt
davor, $n!-1$, ist durch \emph{keine} Zahl von $2$ bis $n$ teilbar.
Der Schritt von $719$ auf $720$ wechselt also mit einem einzigen
Nachfolger von maximaler Teilerfremdheit zu maximaler Teilbarkeit --
und genau an dieser Stelle erzwingt die Horimetrik den Umzug in eine
neue Welt.

\begin{beispiel}
Die Zahl $7$ wohnt minimal in $H_3$ (denn $6\le7<24$) und schreibt
sich dort $013$. In $H_4$ schreibt sie sich $0012$, in $H_5$ als
$00011$, in $H_6$ als $000010$ -- und ab $H_7$ stabil als
$0\cdots07$: Sobald der Horizont groß genug ist, passt die ganze
Zahl in die letzte Ziffer, und weitere führende Nullen sind endlich
wirklich harmlos. Diese Lebensgeschichte erzählen wir später in
einem eigenen Interludium in Ruhe.
\end{beispiel}

Damit steht das Fundament: endliche Welten mit Fakultätsgröße, eine
Bijektion zwischen Ziffern und Werten, und Schalen mit markanten
Grenzen. Was noch fehlt, ist der eigentliche Hauptdarsteller -- die
zweite Identität jeder Zahl. Sie betritt im nächsten Kapitel die
Bühne.

\chapter*{Interludium: Was ein Horizont wirklich ist}
\addcontentsline{toc}{chapter}{Interludium: Was ein Horizont wirklich ist}
\markboth{INTERLUDIUM}{INTERLUDIUM}

\begin{oberflaeche}
Halten wir kurz an, ganz ohne Formeln.

Ein Horizont ist keine Länge und kein Format. Er ist der Raum, in dem
eine Zahl gerade lebt -- mit einer festen, endlichen Anzahl von
Plätzen, die sich aus den Fakultäten ergibt: zwei, sechs,
vierundzwanzig, hundertzwanzig. Eine gewöhnliche Zahl kennt nur
ihren Wert. Eine Hori-Zahl kennt ihren Wert \emph{und} ihren Raum,
und beides gehört zu ihr wie bei einem Buch der Text und das Regal,
in dem es steht.

Das klingt nach einer Kleinigkeit, ist aber ein Perspektivwechsel.
Unsere übliche Mathematik nimmt sich den unendlichen Zahlenstrahl
als gegeben und stellt alle Zahlen gleichzeitig darauf. Die
Horimetrik baut stattdessen Treppen: Jede Stufe ist endlich und
vollständig überschaubar, und die Unendlichkeit steckt nicht in
einer einzelnen Stufe, sondern im Versprechen, dass es immer eine
nächste gibt. Nichts an der klassischen Unendlichkeit ist damit
widerlegt -- sie wird nur anders erzählt: als Folge von Räumen statt
als ein einziger Raum.

Der Rest dieses Buchs untersucht, was dieses Erzählen kostet und was
es einbringt. Es kostet Bequemlichkeit: Beim Umzug von Stufe zu
Stufe ändern alle Ziffern ihre Bedeutung, und man muss sich
entscheiden, was beim Umzug erhalten bleiben soll -- der Wert oder
die Gestalt. Es bringt Struktur ein: Genau aus dieser Entscheidung
entstehen die zwei Identitäten, die zwei Rechenwelten und der Satz,
dass man nicht in beiden gleichzeitig leben kann.
\end{oberflaeche}

\chapter{Zwei Arten zu wachsen: Struktur oder Wert}
\label{ch:erweiterungen}

% Register: D5, S8. Quellen: Spez. Abschnitt 5 und 6.

\section{Der Moment der Gabelung}

Nehmen wir die Sieben in ihrem Zuhause $H_4$: dort schreibt sie sich
$0012$. Jetzt geben wir ihr eine Stelle mehr. Was soll das heißen?

Die erste Antwort ist die naive: Wir stellen eine Null davor.
\[
0012 \;\longmapsto\; 00012 .
\]
Die Ziffern bleiben stehen -- aber in $H_5$ wiegen alle Stellen
außer der letzten anders, und der Wert rutscht von $7$ auf $8$. Die \emph{Gestalt}
wurde bewahrt, der Wert nicht.

Die zweite Antwort ist die vorsichtige: Wir wollen, dass die Zahl
dieselbe bleibt, und codieren die $7$ in $H_5$ neu:
\[
0012 \;\longmapsto\; 00011 .
\]
Der \emph{Wert} wurde bewahrt, die Gestalt nicht.

Im Dezimalsystem sind beide Antworten dieselbe Operation -- deshalb
ist uns nie aufgefallen, dass es zwei Fragen sind. Im Horiraum
trennen sie sich, und wir geben ihnen Namen: Die strukturtreue --
gleichwertig sagen wir: gestalttreue -- Erweiterung $Z$ stellt die
Null voran, der werttreue Lift $L$ codiert um.

\begin{definition}[Register D5]
$Z_n(a_1,\ldots,a_n)=(0,a_1,\ldots,a_n)$ und
$L_n(a)=E_{n+1}(V_n(a))$, wobei $E_{n+1}$ die Codierung in
$\Hn{n+1}$ ist. Es gilt $V_{n+1}(L_n(a))=V_n(a)$, während $Z_n$ im
Allgemeinen den Wert verändert.
\end{definition}

\section{Zwei Treppen, zwei Himmel}

Beide Operationen verbinden die endlichen Horizonte zu einer
unendlichen Treppe -- aber es sind \emph{verschiedene} Treppen, und
sie führen an verschiedene Orte. Die Frage ist jeweils: Welche
Elemente aufeinanderfolgender Horizonte gelten als „dasselbe“?

Unter $L$ gelten zwei Darstellungen als dasselbe, wenn sie denselben
Wert haben: $013$ in $H_3$, $0012$ in $H_4$, $00011$ in $H_5$ sind
ein einziges Objekt -- die Sieben. Unter $Z$ gelten sie als
dasselbe, wenn sie dieselbe Gestalt haben: $12$, $012$, $0012$ sind
ein einziges Objekt -- das Muster „eins-zwei“, dessen Wert von
Stufe zu Stufe wandert.

\begin{satz}[Register S8]\label{satz:grenzen}
Der Grenzraum der Treppe unter $L$ ist $\Nn$, die Menge der
gewöhnlichen natürlichen Zahlen. Der Grenzraum unter $Z$ ist der
Raum aller Gestalten -- formal: der Halbring der Polynome mit
nichtnegativen ganzzahligen Koeffizienten in der Basis der
fallenden Fakultäten, den das nächste Kapitel einführt.
\end{satz}

\begin{proof}[Beweisidee]
Unter $L$ werden genau Darstellungen desselben Wertes identifiziert
(die Codierung ist injektiv), und jeder Wert kommt vor -- die
Klassen \emph{sind} die natürlichen Zahlen. Unter $Z$ werden genau
Ziffernfolgen identifiziert, die sich um führende Nullen
unterscheiden, und jede Gestalt kommt ab ihrem kleinsten zulässigen
Horizont vor -- die Klassen sind die Gestalten. Der vollständige
Beweis über gerichtete Systeme steht in Anhang~\ref{ch:anhang-beweise}.
\end{proof}

\begin{beispiel}
Die beiden Treppen, nebeneinander. Werttreu bleibt die Sieben sie
selbst und wechselt die Gestalt; gestalttreu bleibt das Muster
„eins-zwei“ und wechselt den Wert:
\begin{center}
\begin{tabular}{l|llll}
Welt & $H_2$ & $H_3$ & $H_4$ & $H_5$\\
\hline
werttreu (immer 7) & --- & $013$ & $0012$ & $00011$\\
gestalttreu (immer „12“) & $12=5$ & $012=6$ & $0012=7$ & $00012=8$\\
\end{tabular}
\end{center}
Man beachte die Kreuzung: $0012$ kommt in beiden Zeilen vor -- einmal
als zweite Wohnung der Sieben, einmal als drittes Leben des Musters.
Dasselbe Element, zwei Lesarten.
\end{beispiel}

Das ist der erste Satz dieses Buchs, den man philosophisch lesen
darf: \emph{Dieselben endlichen Räume erzeugen verschiedene
Unendlichkeiten, je nachdem, welchen Konsistenzbegriff man zwischen
den Stufen wählt.} Die vertraute Zahlengerade ist eine dieser
Unendlichkeiten -- die werttreue. Die andere, die gestalttreue, ist
mathematisch genauso legitim, nur hat sie im gewöhnlichen
Zahlensystem nie jemand von der ersten unterscheiden können, weil
dort beide Treppen zusammenfallen.

\section{Die dritte Treppe}

Lange dachten wir, es gäbe genau zwei Arten zu wachsen. Dann las ein
externer Mathematiker eine Kurzbeschreibung dieses Systems -- und
verankerte die Basenleiter versehentlich am anderen Ende. Aus seinem
„Missverständnis“ fiel eine dritte Treppe heraus.

Man kann einer Zahl nämlich auch \emph{rechts} eine Null anhängen.
Dabei behalten alle Ziffern ihre Position von links gezählt, und es
gilt eine verblüffend saubere Regel: Der Wert vervielfacht sich
nicht unkontrolliert, sondern skaliert \emph{exakt} mit dem Faktor $n+2$ --
aus $12$ in $H_2$ (Wert $5$) wird $120$ in $H_3$ (Wert $20$).
Invariant bleibt dabei weder der Wert noch die Gestalt, sondern eine
dritte Größe: der \emph{Fakultätsbruch}
$F=\sum_i a_i/(i+1)!$, denn es gilt stets
\[
V_n(a)=(n+1)!\cdot F(a).
\]
Jede Hori-Zahl ist also auch ein Bruch zwischen null und eins, mal
der Größe ihrer Welt -- und dieses Bruchsystem ist ein alter
Bekannter: Es ist Cantors Fakultätsbasis für Brüche, gut
anderthalb Jahrhunderte alt.

Damit hat der Horiraum drei Treppen, drei Invarianten, drei Himmel:

\begin{center}
\begin{tabular}{l|lll}
Erweiterung & invariant & Wert & Grenzraum\\
\hline
werttreu ($L$) & Wert & $\times1$ & natürliche Zahlen\\
gestalttreu ($Z$, links) & Polynom & ungleichmäßig & Strukturraum\\
bruchtreu ($R$, rechts) & Bruch $F$ & $\times(n+2)$ & Cantor-Brüche\\
\end{tabular}
\end{center}

Und der Einwand „das ist doch seit Cantor bekannt“ bekommt damit
eine präzise Antwort: Er trifft die dritte Treppe vollständig -- und
die anderen beiden samt allem, was dieses Buch über ihr Zusammenspiel
beweist, gar nicht.

\section{Der Eigenhorizont: wo die Gabelung heilt}

Für jede einzelne Zahl gibt es einen Punkt, an dem der Unterschied
zwischen $Z$ und $L$ verschwindet. Steht der gesamte Wert in der
letzten Ziffer -- die Zahl $n$ als $0\cdots0n$; wir nennen diese
Gestalt ihre \emph{Terminaldarstellung}, und sie ist ab dem
Horizont $H_n$ möglich --, dann ist die Nullerweiterung wertneutral:
Gestalt und Wert wachsen ab dort gemeinsam. Wir nennen $H_n$ den
\emph{Eigenhorizont} von $n$. Unterhalb davon wird die Zahl bei
jedem Umzug umgebaut; ab dort wächst sie nur noch still nach links.
Man kann exakt ausrechnen, wann eine führende Null wertneutral ist:
genau dann, wenn alle Ziffern außer der letzten null sind -- und
das nächste Kapitel macht daraus mit der Horizontableitung einen
Einzeiler.

Wie sich das für eine konkrete Zahl anfühlt, erzählt das
Interludium „Die Lebensgeschichte der Sieben“.

\chapter{Strukturpolynome und die Horizontableitung}
\label{ch:strukturpolynome}

% Register: D3, D4, D7, D10, S2, S3, S4, S7, S39, S40, S41, S42, S43.
% Quellen: Spez. 2, 3, 4, 9, Anhänge H und I.

\section{Die Gestalt bekommt Koordinaten}

Bisher war die „Gestalt“ einer Hori-Zahl nur ihr Ziffernbild.
Jetzt geben wir ihr eine algebraische Form -- und damit wird aus
einer Anschauung ein Rechenwerkzeug.

Der Schlüssel sind die \emph{fallenden Fakultäten}
\[
\ff X0=1,\qquad \ff Xd=X(X-1)\cdots(X-d+1),
\]
die Polynome, die „geordnetes Auswählen ohne Zurücklegen“ zählen.
Liest man die Ziffern einer Hori-Zahl von rechts nach links als
Koeffizienten dieser Basis, entsteht ihr \emph{Strukturpolynom}:
\[
a=(a_1,\ldots,a_n)
\quad\longmapsto\quad
P_a(X)=\sum_{d=0}^{n-1}a_{n-d}\,\ff Xd .
\]

\begin{satz}[Auswertungssatz; Register S2]
$V_n(a)=P_a(n+1)$.
\end{satz}

\begin{proof}
Für $d=n-i$ ist $\ff{(n+1)}d=(n+1)!/(i+1)!$ -- genau das Gewicht
der Stelle $i$. Einsetzen liefert die Wertformel.
\end{proof}

Das ist eine kleine Rechnung mit einer großen Konsequenz: Der Wert
einer Hori-Zahl ist \emph{die Auswertung ihrer Gestalt an der Stelle
ihres Horizonts}. Die Sieben in $H_4$, also $0012$, hat die Gestalt
$P(X)=X+2$; ihr Wert ist $P(5)=7$. Stellt man eine Null voran,
bleibt die Gestalt $X+2$ -- nur der Auswertungspunkt wandert:
$P(6)=8$. Die mysteriöse Wertverschiebung durch führende Nullen ist
kein Rätsel mehr, sondern schlicht: \emph{dasselbe Polynom, eine
Stelle weiter ausgewertet.}

Zwei Vorsichtsmaßnahmen. Erstens: Zulässig sind nur nichtnegative
ganzzahlige Koeffizienten, und im Horizont $n$ nur Polynome mit
$d+c_d\le n$ an allen besetzten Stellen ($c_d>0$) --
das ist die Ziffernschranke in Polynomsprache; die Kennzahl
$\sigma(P)=\max_{c_d>0}(d+c_d)$ -- das Maximum läuft über die
besetzten Stellen, und für das Nullpolynom vereinbaren wir
$\sigma(0)=0$ -- heißt Strukturhorizont (Register D4).
Zweitens: Das Polynom allein ist \emph{nicht} die ganze Hori-Zahl.
$12$, $012$ und $0012$ haben dieselbe Gestalt $X+2$, sind aber drei
verschiedene Elemente -- erst das Paar aus Horizont und Polynom
trägt alles.

\section{Die Horizontableitung}

Wie stark verändert eine führende Null den Wert? Die Antwort ist
die diskrete Ableitung, in der Analysis als Vorwärtsdifferenz
bekannt:
\[
\Delta P(X)=P(X+1)-P(X),
\qquad
\Delta\,\ff Xd=d\,\ff X{d-1}.
\]
Die zweite Formel ist das diskrete Ebenbild von
$(x^d)'=d\,x^{d-1}$ -- die fallenden Fakultäten sind für das
Differenzenrechnen, was die Potenzen fürs Differenzieren sind.

\begin{satz}[Register S7]
Die Wertänderung bei strukturtreuer Erweiterung ist
$\Delta P(n+1)$. Insbesondere ist eine führende Null genau dann
wertneutral, wenn $\Delta P=0$, also $P$ konstant ist -- das sind
exakt die Zahlen in oder über ihrem Eigenhorizont.
\end{satz}

\begin{beispiel}
Die Gestalt $X+2$ (Ziffern $\ldots12$) hat Ableitung $1$: Ihre Werte
klettern beim Umziehen um genau eins, $5,6,7,8,\ldots$ Die Gestalt
$\ff X2$ -- das ist die Zwölf in $H_3$, Ziffern $100$ -- hat
Ableitung $2X$: Ihre Werte springen $12, 20, 30, 42,\ldots$, mit
Zuwächsen $8, 10, 12$, die selbst wachsen. Eine Konstante wie die
Gestalt $7$ hat Ableitung $0$ und rührt sich nie.
\end{beispiel}

Damit ist der Eigenhorizont aus Kapitel~\ref{ch:erweiterungen} kein
anschaulicher Begriff mehr, sondern eine Gleichung: Die
\emph{Gestalten} mit Horizontableitung null sind genau die
konstanten Polynome -- als Gestalten gelesen also die natürlichen
Zahlen. (Genau genommen: Als Elemente bilden die Konstanten viele
horizontale Vertreter desselben Werts; erst im gestalttreuen
Grenzraum aus Kapitel~\ref{ch:erweiterungen} wird daraus exakt eine
Kopie von $\Nn$.) Alles andere im Horiraum hat eine
echte „Geschwindigkeit“: Eine Gestalt vom Grad 1 wächst pro
Horizontschritt um einen konstanten Betrag, Grad 2 beschleunigt,
und so fort. Der Polynomgrad ist die \emph{Hori-Dimension} einer
Gestalt (Register D7).

\section{Strecke mal Strecke ist Fläche -- jetzt exakt}

Die Dimension verhält sich unter Multiplikation additiv:
$\dim(PQ)=\dim P+\dim Q$ für $P,Q\ne0$. Was in der Schule eine Eselsbrücke ist --
Länge mal Länge gibt Fläche --, ist hier ein Satz über Grade. Und
man kann direkt mit Gestalten rechnen. Die Vier hat in $H_3$ die
Gestalt $X$ (Ziffern $010$), die Fünf die Gestalt $X+1$ (Ziffern
$011$). Ihr Gestaltprodukt ist
\[
X(X+1)=\ff X2+2X
\]
(denn $X^2=\ff X2+X$), also die Ziffernfolge $120$ -- und
tatsächlich: $120$ in $H_3$ hat den Wert $1\cdot12+2\cdot4+0=20$.
Die Multiplikation $4\cdot5=20$ ist hier ohne einen einzigen
Übertrag geschehen, rein durch Kombinieren der Gestalten. Die
Produktformel der fallenden Fakultäten hat dabei eine hübsche
Nebenbedeutung: $X\cdot X=\ff X2+X$ zerlegt alle geordneten Paare
in die mit verschiedenen und die mit gleichen Einträgen.

\section{Was Operationen kosten: das Horizontkalkül}
\label{sec:horizontkalkuel}

Jede Gestalt braucht einen Horizont, der groß genug ist -- die
Kennzahl dafür ist der Strukturhorizont $\sigma$. Damit bekommt jede
Operation einen \emph{Preis}: Wie viel Kapazität muss der Zielraum
haben, damit das Ergebnis dort sicher Platz findet? Für einen
Operator $T$ fragen wir nach dem schlimmsten Fall
\[
\Gamma_T(r,s)=\max\{\sigma(T(P,Q)) \mid \sigma(P)\le r,\ \sigma(Q)\le s\}
\]
(Register D10; für einstellige Operatoren sinngemäß $\Gamma_T(r)$). Auf den ersten Blick ein Ungetüm: Schon für $r=s=8$
wären über hundert Milliarden Paare zu prüfen. Es geht mit einem
einzigen.

\begin{satz}[Extremalprinzip; Register S39]
In jedem Horizont $r$ gibt es die \emph{maximal gefüllte} Gestalt
$M_r$ -- jede Ziffer auf Anschlag,
$M_r(X)=\sum_{d<r}(r{-}d)\,\ff Xd$; wir nennen sie auch das
\emph{Maximalpolynom} des Horizonts.
Für jeden Operator, der auf wachsende Koeffizienten mit wachsenden
Koeffizienten antwortet (das tun alle aus Addition, Multiplikation
und $\Delta$ zusammengesetzten), gilt
\[
\Gamma_T(r,s)=\sigma\bigl(T(M_r,M_s)\bigr).
\]
\end{satz}

Der schlimmste Fall unter Abermilliarden Kandidaten ist immer
derselbe: alle Ziffern voll. Der Beweis (Anhang~\ref{ch:anhang-beweise})
ist eine Zeile Monotonie -- aber die Konsequenz ist ein
\emph{Kalkül}: Addition kostet exakt $r+s$; die Multiplikation kostet
$\beta(r,s)$, die im Kern fakultätisch wachsende Produkthöhe, die uns
in Kapitel~\ref{ch:zaehlungen} als Zählfolge wiederbegegnet. Und die
Ableitung liefert die größte Überraschung des Kalküls:

\begin{satz}[Differenzüberlauf; Register S40]
Für $r\ge2$ gilt
\[
\max_{\sigma(P)\le r}\sigma(\Delta P)
=\Bigl\lfloor\tfrac{(r+1)^2}{4}\Bigr\rfloor-1 .
\]
\end{satz}

Man halte kurz inne: $\Delta$ \emph{senkt} den Grad -- und lässt die
Kapazität trotzdem quadratisch anschwellen. Schon bei $r=8$ hebt eine
einzige Differenz den Strukturhorizont von 8 auf 19. Der Zeuge ist
eine Gestalt mit einer einzigen Ziffer in der Mitte: Die Ableitung
multipliziert jeden Koeffizienten mit seinem Grad, und am teuersten
ist der Kompromiss aus „hoher Grad“ und „großer Koeffizient“ genau
in der Mitte -- daher der Term $\lfloor(r{+}1)^2/4\rfloor$, das
größte Produkt zweier Zahlen mit fester Summe. Die Lehre: Der Grad
misst die \emph{Dimension} einer Gestalt, $\sigma$ misst ihre
\emph{Kapazität}, und die beiden können gegenläufig springen. Eine
gewöhnliche Gradbetrachtung ist für diesen Effekt blind.

Eine Operation fehlt dem Kalkül noch, die heikelste: das
\emph{Einsetzen} einer Gestalt in eine andere, $P\circ Q$. Dass das
überhaupt erlaubt ist, ist selbst ein Satz -- es ist von vornherein
nicht klar, dass beim Einsetzen wieder lauter zulässige, ganzzahlige,
nichtnegative Ziffern entstehen.

\begin{satz}[Kompositionsabschluss; Register S41]
Setzt man eine Gestalt in eine Gestalt ein, entsteht wieder eine
Gestalt. Auch das Einsetzen gehorcht dem Extremalprinzip:
$\Gamma_\circ(r,s)=\sigma(M_r\circ M_s)$.
\end{satz}

Der Beweis (Anhang~\ref{ch:anhang-beweise}) ist das kombinatorische
Herzstück dieses Kapitels: Man liest $Q$ als Bauplan-Vorrat --
jede Ziffer $a_k$ stellt $a_k$ Schablonen bereit, die auf $k$
verschiedene Plätze montiert werden --, dann zählt $P\circ Q$
Tupel solcher Montagen, und die nötige Teilbarkeit durch $m!$
kommt daher, dass das Umbenennen der Plätze keine Montage-Auswahl
festhalten kann. Die Kosten des Einsetzens sind dann allerdings von
ganz anderem Kaliber als alles Bisherige: Die Diagonale
$\Gamma_\circ(r,r)$ beginnt mit
\[
1,\quad 4,\quad 23,\quad 6541,\quad 477\,149\,751,\quad\ldots
\]
-- Einsetzen multipliziert die Grade \emph{und} lässt die
Koeffizienten explodieren; das ist die neunte und mit Abstand am
schnellsten wachsende Zählfolge dieses Projekts
(Kapitel~\ref{ch:zaehlungen}). Der harmloseste Spezialfall, das
bloße Verschieben $P(X)\mapsto P(X{+}1)$, kostet exakt
$r+\lfloor r^2/4\rfloor$ -- auf den Punkt so viel wie eine Ableitung
einen Horizont höher (Register S42), eine kleine Dualität, die das
Kalkül gratis abwirft.

\section{Was die Gestalten zählen: die Zählbreite}
\label{sec:zaehlbreite}

Zum Abschluss des Kapitels die Frage, die man einem algebraischen
Objekt immer stellen sollte: Zählt es etwas? Die Antwort kam als
externe Einordnung ins Projekt und ist die bisher beste Brücke der
Horimetrik in ein etabliertes Gebiet -- die \emph{Zählkomplexität}.

Der Schlüssel ist die Grundbedeutung der fallenden Fakultät:
$n^{\underline d}$ zählt die geordneten $d$-Tupel
\emph{verschiedener} Objekte aus einem Vorrat von $n$. Ist nun $n$
die Anzahl der Lösungen irgendeines Suchproblems -- in der
Informatik sagt man: der \emph{Zeugen} --, dann liest sich eine
Gestalt $P=\sum_d c_d\ff Xd$ als Bauplan für neue Zählobjekte:
$c_0$ feste Sonderobjekte, $c_1$ Farben für jede einzelne Lösung,
$c_2$ Varianten für jedes geordnete Paar verschiedener Lösungen,
und so fort. $P(n)$ ist exakt die Anzahl dieser gefärbten
Zeugen-Tupel. Diese harmlose Beobachtung hat eine präzise
Konsequenz:

\begin{satz}[Hori-\#P-Satz; Register S43]
Ist $f$ die Zeugen-Zählfunktion eines beliebigen NP-Problems und
$P$ eine Gestalt, so ist auch $P\circ f$ die Zeugen-Zählfunktion
eines NP-Problems. Jede Gestalt ist also ein
\emph{Abschlussoperator} der Zählklasse \#P.
\end{satz}

Der Beweis (Anhang~\ref{ch:anhang-beweise}) ist genau die
Bauplan-Lektüre: rate den Grad, rate die Farbe, rate ein Tupel
verschiedener Zeugen. Ehrlichkeit zuerst: Die \emph{qualitative}
Seite dieser Aussage ist bekannt -- Polynome mit nichtnegativen
ganzzahligen Koeffizienten in der \emph{Binomial}basis heißen in
der Literatur \emph{binomial-good} und sind als die
(relativierenden) \#P-Abschlussoperationen charakterisiert. Wegen
$\ff Xd=d!\binom Xd$ sind unsere Gestalten eine \emph{echte
Teilklasse} davon: Die Teilbarkeit durch $d!$ ist genau der
Unterschied zwischen geordneten Tupeln und ungeordneten Auswahlen.
Was die Literatur offenbar \emph{nicht} hat, ist die quantitative
Schicht -- und die liefert die Horimetrik gratis mit:

Plötzlich bedeuten alle Größen dieses Kapitels etwas. Die
\emph{Addition} von Gestalten ist die disjunkte Vereinigung zweier
Zählkonstruktionen. Die \emph{Multiplikation} koppelt zwei
Konstruktionen, und ihre Produktformel-Koeffizienten zählen exakt
die Überlappungsmuster zweier Zeugengruppen -- das
$\beta$-Wachstum misst also, wie die Überlappungsvielfalt beim
Koppeln explodiert. Die \emph{Horizontableitung} zählt, welche
neuen Strukturen ein einziger zusätzlicher Zeuge ermöglicht
($d$ Positionen mal Rest-Tupel: genau $d\,\ff n{d-1}$) -- und der
Differenzüberlauf besagt: Ein neuer Zeuge senkt zwar die
Tupelordnung, kann die Variantenvielfalt aber quadratisch aufblähen.
Und $\sigma$ selbst? Zeichnet man das Koeffizientenbrett -- Spalte
$d$ mit $c_d$ Kästchen --, so ist $\sigma(P)=\max(d+c_d)$ die
kleinste \emph{Treppe}, in die das Brett passt: das Maximum aus
„wie viele Zeugen gleichzeitig“ und „wie viele Varianten davon“.
Wir nennen diese Größe die \emph{Zählbreite} einer Konstruktion.
Jeder Horizont $r$ ist damit eine endliche, vollständig
aufzählbare Familie von genau $(r{+}1)!$
\#P-Abschlussoperatoren, und das Horizontkalkül aus
Abschnitt~\ref{sec:horizontkalkuel} beziffert optimal, wie die
Zählbreite beim Vereinigen, Koppeln, Erweitern und Einsetzen
wächst.

Zwei Sätze zur Abgrenzung, damit kein falscher Zauber entsteht: In
dieser Nachbarschaft wohnt eine gewichtige offene Vermutung (schon
ihre einfachste Fassung würde P$\,\ne\,$NP implizieren), und
unsere Operatoren liegen ausdrücklich auf der \emph{gut
verstandenen} Seite dieser Landschaft -- die Horimetrik löst hier
gar nichts. Was sie beiträgt, ist bescheidener und vielleicht
gerade deshalb brauchbar: aus einem Erlaubt-oder-nicht-erlaubt
erstmals ein \emph{Wie breit} zu machen.

\begin{oberflaeche}
Was dieses Kapitel gebracht hat, in einem Satz: Die Gestalt einer
Hori-Zahl ist ein Polynom, ihr Wert ist dieses Polynom am Ort ihres
Horizonts ausgewertet, und die rätselhafte Wirkung führender Nullen
ist einfach die Ableitung dieses Polynoms. Die gewöhnlichen Zahlen
sind die Gestalten mit Ableitung null -- die einzigen Bewohner des
Horiraums, denen es egal ist, wo sie wohnen. Alle anderen bewegen
sich, wenn ihre Welt wächst. Jede Rechenoperation hat einen exakt
bezifferbaren Preis in Kapazität, den stets dasselbe extreme
Objekt bezahlt: die randvolle Gestalt. Und zuletzt hat sich
gezeigt, dass die Gestalten etwas \emph{zählen} -- gefärbte
Gruppen verschiedener Lösungen eines Suchproblems -- und $\sigma$
die Breite dieser Zählkonstruktionen misst. Im nächsten Kapitel
lassen wir diese bewegten Objekte rechnen -- und stoßen auf die
Frage, was dabei mit ihrer Vergangenheit geschehen soll.
\end{oberflaeche}

\chapter*{Interludium: Die Lebensgeschichte der Sieben}
\addcontentsline{toc}{chapter}{Interludium: Die Lebensgeschichte der Sieben}
\markboth{INTERLUDIUM}{INTERLUDIUM}

\begin{oberflaeche}
Begleiten wir eine einzige Zahl durch alle ihre Wohnungen.

Die Sieben kommt in der Welt mit vierundzwanzig Plätzen zur Welt --
früher geht es nicht, die Welt davor endet bei fünf. Dort
heißt sie $013$. Es ist eine beengte Wohnung: Drei Ziffern, alle
tragen etwas bei, nichts ist Dekoration.

Dann zieht sie um, Welt für Welt, immer werttreu -- sie will ja
dieselbe Sieben bleiben. Und bei jedem Umzug muss sie sich neu
einrichten: In der Welt mit hundertzwanzig Plätzen heißt sie
$0012$, in der nächsten $00011$, dann $000010$. Wer die Umzüge
nebeneinanderlegt, sieht ein rastloses Muster -- die Ziffern wandern,
tauschen, schrumpfen. Kein Umzug ist wie der andere, und nichts an
der Schreibweise verrät auf den ersten Blick, dass es immer dieselbe
Zahl ist.

Bis zur Welt mit vierzigtausenddreihundertzwanzig Plätzen. Dort
passiert etwas Neues: Die Sieben passt zum ersten Mal vollständig in
die letzte Ziffer -- sie heißt jetzt $0000007$, und man sieht ihr
ihren Wert wieder an. Von diesem Umzug an ändert sich nie mehr
etwas. Jede weitere Welt hängt nur noch eine stille Null vorne an:
$00000007$, $000000007$, immer so weiter. Die Sieben ist
angekommen; wir sagen: Sie hat ihren Eigenhorizont erreicht.

Jede Zahl hat diese Biographie: eine rastlose Jugend, in der jeder
größere Raum eine neue Identitätsverhandlung bedeutet, und ein
ruhiges Erwachsenenalter, ab dem mehr Raum einfach nur mehr Raum
ist. Große Zahlen haben eine lange Jugend -- die Zahl
siebenhundertzwanzig zum Beispiel wird bis zur Welt Nummer
siebenhundertzwanzig
umgebaut, und auf diesem langen Weg trägt sie zeitweise Gestalten,
die ihr niemand zutraut: In der neunten Welt ist sie schlicht eine
einzelne Eins an der vierten Stelle von rechts, denn zehn mal neun
mal acht ist siebenhundertzwanzig. Genau diese Doppelnatur -- dem
Wert nach sperrig, der Gestalt nach schlicht -- hat sich später als
Schlüssel zu einem der Hauptsätze erwiesen.

Merken muss man sich aus alledem nur eines: Im Horiraum ist
„dieselbe Zahl“ keine Selbstverständlichkeit, sondern eine
Entscheidung. Man kann den Wert festhalten und die Gestalt opfern,
oder umgekehrt. Wie viel diese Entscheidung kostet, ist keine
Geschmacksfrage mehr, sondern ein Satz -- aber das gehört ins
nächste Kapitel.
\end{oberflaeche}

\chapter{Arithmetik mit Gedächtnis: der Halbring}
\label{ch:arithmetik}

% Register: D6, S5, S6, S10-S16, V9. Quellen: Spez. 8, Anhang A-C.

Bis hierher haben wir Hori-Zahlen nur betrachtet und verschoben.
Jetzt wollen wir mit ihnen rechnen. Das klingt harmlos,
ist aber die Stelle, an der die Horimetrik ihre erste echte
Entscheidung verlangt: Wenn zwei Hori-Zahlen addiert werden, welchen
Horizont bekommt das Ergebnis?

Im gewöhnlichen Zahlensystem stellt sich diese Frage nicht -- dort hat
ein Ergebnis keinen Horizont, nur einen Wert. Hier dagegen ist der
Horizont Teil der Zahl. Eine Rechenvorschrift muss deshalb aus zwei
Teilen bestehen: einer Regel für den Wert und einer Regel für den
Horizont. Die Horizontregel nennen wir einen \emph{Selektor}.

\section{Der erste Fehlschlag und was er lehrt}

Der naheliegendste Selektor ist der konservative: Das Ergebnis erbt
den größten beteiligten Horizont, notfalls vergrößert um das, was der
Ergebniswert mindestens braucht. Für die Addition funktioniert das
tadellos. Für die Multiplikation scheitert es -- und zwar an der
unschuldigsten aller Zahlen.

\begin{sackgasse}
Multipliziert man konservativ, dann gilt im Allgemeinen
$(a\htimes b)\htimes c\ne a\htimes(b\htimes c)$, sobald eine Null
beteiligt ist. Die Null löscht den \emph{Wert} jedes Produkts, aber
der konservative Selektor schleppt die \emph{Horizonte} der
Zwischenschritte weiter mit. Je nachdem, wann die Null zuschlägt,
erinnert sich das Ergebnis an verschiedene Rechenwege -- gleicher
Wert, verschiedene Horizonte, keine Klammerfreiheit. Lehre: Ein
Horizont trägt Rechengeschichte, und Rechengeschichte verträgt sich
nicht automatisch mit den Rechengesetzen.
\end{sackgasse}

\begin{beispiel}
Konkret, mit der konservativen Multiplikation (Horizont = Maximum
der Beteiligten, notfalls vergrößert): Es sei $a=b=(2,5)$ -- die
Fünf mit Horizont 2 -- und $c=(1,0)$ die Null mit Horizont 1. Links
geklammert: $a\htimes b=(4,25)$, denn $25$ braucht Horizont 4;
danach $(4,25)\htimes c=(4,0)$. Rechts geklammert:
$b\htimes c=(2,0)$, danach $a\htimes(2,0)=(2,0)$. Beide Wege enden
beim Wert null -- aber einmal mit Horizont 4, einmal mit Horizont 2.
$(4,0)\ne(2,0)$: keine Assoziativität.
\end{beispiel}

\section{Der wertseitige Halbring}

Der Ausweg ist eine Arbeitsteilung: Die Addition darf sich erinnern,
die Multiplikation muss vergessen. (Eine Konvention vorweg, die für
das ganze Buch gilt: \emph{Gesetzestreu} nennen wir eine Arithmetik,
wenn beide Operationen assoziativ und kommutativ sind, die
Multiplikation über die Addition distribuiert und ein additiv
neutrales Element existiert. Eine multiplikative Eins und das
Schlucken durch die additive Null verlangen wir \emph{nicht}
grundsätzlich, sondern nur dort, wo sie ausdrücklich genannt
werden -- beides wird uns gleich als echte Unterscheidungsmerkmale
begegnen.)

\begin{satz}[wertseitiger Halbring]\label{satz:a3}
Auf dem Horiraum $\Hori=\{(n,x)\mid n\ge\muh(x)\}$ seien, in
Wertkoordinaten $(n,x)$,
\[
(n,x)\hplus(m,y)=\bigl(\max(n,m,\muh(x+y)),\,x+y\bigr),
\qquad
(n,x)\htimes(m,y)=\bigl(\muh(xy),\,xy\bigr).
\]
Dann sind $\hplus$ und $\htimes$ kommutativ und assoziativ,
$\htimes$ distribuiert über $\hplus$, $(0,0)$ ist neutral für
$\hplus$ und absorbierend für $\htimes$.
\end{satz}

Der Beweis ist erstaunlich kurz; sein einziger Kern ist die
Monotonie des Minimalhorizonts: mehr Wert braucht nie weniger Platz
(vollständig in Anhang~\ref{ch:anhang-beweise}).
Bemerkenswert ist, was diesem Halbring \emph{fehlt}: eine Eins. Es
gibt kein Element, das bei Multiplikation alles unverändert lässt.
Stattdessen gilt für die Hori-Zahl Eins -- also $(1,1)$, die Eins in
ihrem kleinsten Zuhause -- die Gleichung
\[
h\htimes(1,1)=\KV(h)
\qquad\text{mit}\quad
\KV(n,x)=(\muh(x),\,x).
\]
Multiplikation mit der Eins ist die
Wertkompaktifizierung.\footnote{Terminologische Warnung:
„Kompaktifizierung“ ist hier nicht topologisch gemeint. Formal
handelt es sich um idempotente Retraktionen auf den jeweiligen
Minimalhorizont -- wir behalten das anschauliche Wort bei und bitten
topologisch geschulte Leser um Nachsicht.} Die Eins
wirkt nicht als neutrales Element, sondern als \emph{Projektor}: Sie
drückt jede Hori-Zahl auf ihre kompakte Form. Die kompakten Elemente
bilden darunter einen eigenen kleinen Halbring, und dieser ist exakt
das gewöhnliche $\Nn$. Die vertrauten natürlichen Zahlen sitzen also
als „Eck“ im Horiraum, und der Weg dorthin ist keine Zusatzoperation,
sondern Multiplikation mit der Eins.

\begin{beispiel}
Rechnen wir einmal durch. $(3,7)\hplus(2,5)$: Wert $12$, Horizont
$\max(3,2,\muh(12))=\max(3,2,3)=3$, also $(3,12)$ -- der größte
beteiligte Horizont überlebt. $(3,7)\htimes(2,5)$: Wert $35$,
Horizont $\muh(35)=4$, also $(4,35)$ -- die Vergangenheit ist
vergessen. Und der Projektor: Die Sieben mit unnötig großem
Horizont, $(5,7)$, mal die Eins $(1,1)$ ergibt
$(\muh(7),7)=(3,7)$ -- die Multiplikation mit Eins hat aufgeräumt.
\end{beispiel}

\begin{oberflaeche}
Was ist hier passiert? Wir haben eine Arithmetik gefunden, in der
Zahlen sich beim Addieren an ihre Vergangenheit erinnern -- der
größte beteiligte Horizont überlebt -- und beim Multiplizieren neu
geboren werden. Und das gewöhnliche Rechnen, das wir seit der
Grundschule kennen, ist darin enthalten: als der Teil des Horiraums,
der keine Vergangenheit mit sich trägt.
\end{oberflaeche}

\section{Exzess: das Gedächtnis bekommt eine Koordinate}

Lange sah es so aus, als sei diese vergessende Multiplikation
alternativlos. Die Vermutung lag nahe, dass die Rechengesetze jede
andere Horizontregel zerstören. Sie ist falsch, und die Widerlegung
ist lehrreicher als ein Beweis es gewesen wäre.

Der Horizont jeder Hori-Zahl zerfällt kanonisch in zwei Anteile:
\[
n=\underbrace{\muh(x)}_{\text{Pflicht}}
 +\underbrace{\varepsilon}_{\text{Exzess}} .
\]
Der Pflichtanteil ist das Minimum, das der Wert braucht; der Exzess
ist freiwillig mitgeführtes Gedächtnis. Diese Zerlegung ist keine
Konvention, sondern durch die Fakultätsschalen eindeutig bestimmt --
und sie entkoppelt die Frage nach dem Selektor vollständig: Man darf
auf dem Exzess \emph{eigenständig} rechnen.

\begin{satz}[Produktprinzip]\label{satz:produktprinzip}
Jede kommutative Halbringstruktur auf den Exzessen liefert eine
gesetzestreue Hori-Arithmetik: Werte rechnen gewöhnlich,
Exzesse rechnen in der gewählten Struktur.
\end{satz}

Das Produktprinzip liefert damit eine ganze Familie neuer
Arithmetiken. Zusammen mit dem schon konstruierten vergessenden
Halbring aus Satz~\ref{satz:a3} -- der, wohlgemerkt, \emph{keine}
Produktkonstruktion ist: Sein Additions-Selektor vergleicht die
Horizonte selbst, nicht nur die Exzesse -- ergeben sich drei
besonders natürliche Arithmetiken:

\begin{center}
\begin{tabular}{lccc}
 & Gedächtnis bei $+$ & Gedächtnis bei $\times$ & Eins \\
\hline
vergessend (Satz~\ref{satz:a3}; kein Produkt) & Maximum (der Horizonte) & stirbt & keine \\
bewahrend (tropisch) & Maximum & addiert sich & $(1,1)$ \\
verstärkend (arithmetisch) & addiert sich & multipliziert sich & $(2,1)$ \\
\end{tabular}
\end{center}

\begin{beispiel}
Es sei $a=(5,7)$ -- Pflicht $\muh(7)=3$, Exzess $2$ -- und
$b=(3,5)$ -- Pflicht $2$, Exzess $1$. Das Produkt hat den Wert $35$
mit Pflicht $\muh(35)=4$. Die drei Arithmetiken unterscheiden sich
nur im Exzess des Ergebnisses: vergessend $(4,35)$ (Exzess $0$),
bewahrend $(4{+}2{+}1,35)=(7,35)$ (Exzesse addieren sich),
verstärkend $(4{+}2,35)=(6,35)$ (Exzesse multiplizieren sich:
$2\cdot1=2$).
\end{beispiel}

Die bewahrende Arithmetik hat eine Kuriosität: Ihre Null schluckt
nicht. Multipliziert man eine Hori-Zahl mit null, wird der Wert
ausgelöscht, aber das Gedächtnis überlebt. Die verstärkende dagegen
ist ein vollwertiger kommutativer Halbring mit Eins und schluckender
Null -- und ihre Multiplikation löscht das Gedächtnis nicht
pauschal, sondern verrechnet es arithmetisch. Ganz gedächtnistreu
ist freilich auch sie nicht: Ein Operand mit Exzess null absorbiert
das Gedächtnis des anderen, denn null mal irgendetwas bleibt null --
genau daher rührt ihre schluckende Null.

\begin{oberflaeche}
Die Frage „welchen Horizont bekommt das Ergebnis?“ hat sich in
etwas viel Schöneres verwandelt: „Was soll mit der Vergangenheit
geschehen, wenn gerechnet wird?“ Und die Antwort der Horimetrik ist:
Das darfst du wählen -- vergessen, bewahren oder verstärken -- und
jede Wahl ergibt eine in sich vollkommen widerspruchsfreie Welt.
\end{oberflaeche}

\section{Die Dualität: dieselbe Eins, zwei Projektoren}

Bis hierher rechnete alles auf der Wertseite. Aber Hori-Zahlen
tragen neben dem Wert auch ihre Gestalt (und, seit
Kapitel~\ref{ch:erweiterungen}, den Bruch -- von dem in
Abschnitt~\ref{sec:seitentreue} noch die Rede sein wird) -- und die Gestaltseite hat
ihre eigene Arithmetik. Man
kann statt der Werte die \emph{Strukturpolynome} addieren und
multiplizieren; das ist das ziffernweise Rechnen ohne Überträge, bei
dem etwa aus den Strukturen von $4$ und $5$ direkt die Struktur von
$20$ entsteht.

\begin{satz}[Dualitätssatz]\label{satz:dualitaet}
Der Horiraum trägt einen zweiten, strukturseitigen Halbring: Addition
konservativ, Multiplikation kompakt -- aber gemessen am
Strukturhorizont $\sigmah$ statt am Werthorizont $\muh$. Gerechnet
wird in Strukturkoordinaten $(n,P)$; über die Ziffern-Bijektion aus
Kapitel~\ref{ch:horizonte} bezeichnen $(n,x)$ und das Paar aus $n$
und der Gestalt von $x$ dasselbe Element. In ihm gilt
\[
h\boxtimes(1,1)=\KS(h)
\qquad\text{mit}\quad
\KS(n,P)=(\sigmah(P),\,P),
\]
und $\KS$ ist mit beiden Operationen verträglich. Das Eck dieses
Halbrings ist der Strukturpolynomraum.
\end{satz}

Die Pointe verdient einen eigenen Absatz. Es ist \emph{dieselbe}
Hori-Zahl $(1,1)$ -- die Eins --, die auf der Wertseite den Projektor
$\KV$ und auf der Strukturseite den Projektor $\KS$ trägt. Zwei
Multiplikationen, ein Element. Und das älteste Ergebnis dieses
Projekts -- dass die beiden Aufräum-Operationen nicht vertauschen;
der nächste Abschnitt macht das formal -- entpuppt sich als das
Nichtkommutieren dieser zwei Multiplikationen mit ein und derselben
Eins.

\section{Aufräumen: der Kommutator und das Tiefengesetz}
\label{sec:kommutator}

Weil auf ihnen ein ganzer Themenstrang dieses Buchs ruht, stellen
wir die beiden Projektoren jetzt nebeneinander (Register D6):
\[
\KV(n,x)=(\muh(x),\,x),
\qquad
\KS(n,P)=(\sigmah(P),\,P).
\]
$\KV$ räumt \emph{wertkompakt} auf: gleicher Wert, kleinster
Horizont -- die Gestalt wird dabei neu codiert. $\KS$ räumt
\emph{strukturkompakt} auf: gleiche Gestalt, kleinster Horizont --
der Wert wandert dabei mit dem Auswertungspunkt nach unten. Beide
sind idempotent, aber sie kommutieren nicht, und die Abweichung ist
messbar. Wir nennen
\[
\delta(h)=\hor\bigl(\KS\KV(h)\bigr)-\hor\bigl(\KV\KS(h)\bigr)
\]
den \emph{Kompaktifizierungs-Kommutator}; seine negativen Werte
heißen \emph{Tiefe}.

\begin{beispiel}
Der Hausgeist dieses Buchs: $h=(6,6)$, die Sechs in
Terminaldarstellung. $\KV$ zuerst: $(6,6)\mapsto(3,6)$; dort hat
die Sechs die Gestalt $X{+}2$ mit $\sigmah=2$, also
$\KS\KV(h)=(2,5)$ -- Horizont~$2$. $\KS$ zuerst: Die konstante
Gestalt $6$ ist schon strukturkompakt, $\KS(h)=h$, danach
$\KV(h)=(3,6)$ -- Horizont~$3$. Also $\delta(h)=2-3=-1$: Die
Reihenfolge des Aufräumens entscheidet, wo man landet -- und als
was.
\end{beispiel}

Nach oben wächst $\delta$ vergleichsweise schnell; der kleinste
Zeuge für $+2$ ist die $720$ als Gestalt $\ff X3$ im Horizont $9$
(die „einzelne Eins an der vierten Stelle“ aus dem letzten
Interludium). Nach unten dagegen ist $\delta$ fakultätisch träge,
und diese Trägheit gehorcht einem exakten Gesetz. Nenne
\[
g(m)=m-\min\{s\mid M_s(m{+}1)\ge m!\}
\]
die \emph{Schalenlücke} der Schale $S_m$: wie weit unter $m$ die
Kapazität einer Gestalt liegen darf, deren Wert trotzdem noch $S_m$
erreicht ($M_s$ ist die maximal gefüllte Gestalt aus
Kapitel~\ref{ch:strukturpolynome}).

\begin{satz}[Tiefengesetz, untere Seite; Register S33]
\label{satz:tiefengesetz}
Für jede Schale $m\ge3$ gilt
\[
\min\{\delta(h)\mid \muh(\val(h))=m\}=-g(m),
\]
und ein stets tauglicher Zeuge -- im Allgemeinen nicht der
einzige -- ist die Terminaldarstellung des Maximalpolynom-Wertes
$M_{m-g(m)}(m{+}1)$. Die Sprungstellen $m_g=\min\{m\mid g(m)\ge g\}$
bilden eine Zählfolge mit bewiesener Asymptotik
$m_g/(g{+}2)!\to1$ (Beweise in Anhang~\ref{ch:anhang-beweise}).
\end{satz}

Für Schale $3$ ist der Zeuge genau der Hausgeist von oben, und
$m_1=3$: Ab der dritten Schale ist Tiefe $-1$ möglich, ab der
fünfzehnten $-2$, ab der vierundachtzigsten $-3$. Die
\emph{beidseitigen Fixpunkte} schließlich sind die Elemente, denen
das Aufräumen nichts mehr anhaben kann -- fixiert von $\KV$
\emph{und} $\KS$ zugleich, also $n=\muh(x)=\sigmah(P)$. Beide
Begriffe, Sprungstellen wie Fixpunkte, werden in
Kapitel~\ref{ch:zaehlungen} gezählt.

Erwähnt sei noch das Feinkorn dieser Nicht-Kommutativität: Das von
$\KV$ und $\KS$ erzeugte Transformationsmonoid -- das
\emph{Kompaktifizierungsmonoid} -- ist endlich: genau sechs
Elemente, von denen alle außer $\KV\KS$ idempotent sind
(Register S21; Beweis in Anhang~\ref{ch:anhang-beweise}). Der
Zeuge mit dem kleinsten \emph{Wert} ist wieder der Hausgeist
$(6,6)$; in der kleinsten \emph{Welt} scheitert die Idempotenz von
$\KV\KS$ schon an $(5,25)$.

\section{Seitentreue}
\label{sec:seitentreue}

Kann man die beiden Seiten verweben -- etwa wertseitig addieren und
strukturseitig multiplizieren? Gesetzestreu heißt dabei, wie
eingangs vereinbart: beide Operationen kommutativ und assoziativ,
die Multiplikation distribuiert, eine additive Null existiert. Der
systematische Test aller sechzehn Kombinationen -- Wert oder
Gestalt als Payload, konservativ oder kompakt als Selektor, für
jede der beiden Operationen -- gibt dann eine klare Antwort: Genau
zwei \emph{dieser sechzehn} sind gesetzestreu, der wertseitige
Halbring und sein strukturseitiger Spiegel. (Die Produktfamilie aus
Satz~\ref{satz:produktprinzip} steht außerhalb dieses Rasters --
ihre Selektoren rechnen auf Exzessen -- und ist durchweg
gesetzestreu, bleibt aber ganz auf der Wertseite.) Alle acht gemischten
Kandidaten sterben an der Distributivität; die übrigen sechs
seitenreinen scheitern an der Neutralität (kompakte Addition
vergisst den Horizont des Operanden) oder an der Assoziativität
(konservative Multiplikation, das Nullproblem vom Kapitelanfang).

\begin{beispiel}
Der paradigmatische Zeuge, Ziffer für Ziffer. ($\oplus$ und
$\otimes$ bezeichnen in diesem Beispiel die gemischte
Kandidaten-Arithmetik, nicht die Operationen aus
Satz~\ref{satz:a3}; \emph{wertgetragen} heißt: die Operation
verrechnet die Werte, \emph{strukturgetragen}: die Gestalten.)
Es seien $a=b=(1,1)$ und
$c=(2,5)$, dessen Gestalt $X+2$ ist (Ziffern $12$). Addiert man
wertgetragen, ist $b\oplus c=(3,6)$; die Sechs hat in $H_3$
zufällig ebenfalls die Gestalt $X+2$ (Ziffern $012$). Multipliziert
man nun strukturgetragen mit $a$ (Gestalt $1$), bleibt die Gestalt
$X+2$ und wird kompakt bei Horizont 2 ausgewertet: Ergebnis
$(2,5)$ -- \emph{der Wert ist von 6 auf 5 gefallen}, weil die
Strukturmultiplikation nur die Gestalt kennt. Rechts dagegen:
$a\otimes b=(1,1)$, $a\otimes c=(2,5)$, wertgetragen addiert
$(3,6)$. Links $(2,5)$, rechts $(3,6)$: Die Distributivität ist
tot, und man sieht genau, woran -- die eine Seite rechnet mit dem
Wert ehrlich weiter, die andere hat ihn beim Gestaltwechsel
verloren.
\end{beispiel}

\begin{satz}[Seitentreue-Satz]\label{satz:seitentreue}
Jede gesetzestreue Arithmetik \emph{der hier betrachteten
payload-getra\-genen Form} -- die Operationen tragen Wert, Gestalt
oder Bruch, die Horizontregeln sind beliebig -- ist seitenrein:
Addition und Multiplikation operieren beide auf der Wertseite oder
beide auf der Strukturseite.
\end{satz}

Die Einschränkung auf payload-getragene Operationen ist wesentlich:
Auf einer abzählbaren Menge lässt sich per Transport jede beliebige
Halbringstruktur definieren -- der Satz spricht über die
Arithmetiken, die Wert, Gestalt oder Bruch \emph{respektieren},
und über diese vollständig.

Zur genauen Reichweite: Der Beweis braucht je nach Fall sogar
weniger als die volle Gesetzestreue -- ein Mischfall stirbt an der
Links-Distributivität allein, der andere an der beidseitigen
Distributivität (die zu jedem, auch nichtkommutativen, Halbring
gehört): Die Rechts-Distributivität spiegelt die Kettenkonstruktion
und ersetzt die Kommutativität. Auch nichtkommutative Mischformen
sind damit ausgeschlossen; denkbar bleibt allein ein nur
\emph{einseitig} distributives Mischsystem -- eine Kuriosität
außerhalb jeder Halbring-Axiomatik.

Der Beweis verdient es, in seiner Grundidee erzählt zu werden, denn
er hat eine Pointe (vollständig in Anhang~\ref{ch:anhang-beweise}). Zunächst sah es schlecht aus: Einzelne
Rechnungen, die eine gemischte Arithmetik erfüllen müsste, sind
tatsächlich lösbar -- teils über Pell-Gleichungen, also genau die
Sorte Zahlentheorie, die gerne Ausnahmen produziert. Ein direkter
Widerspruch schien unerreichbar. Die entscheidende Waffe ist
verblüffend schlicht: die \emph{Konstanten}. Multipliziert man eine
per Addition aufgebaute Kette von Zahlen mit immer größeren
konstanten Strukturen, dann treibt die Ziffernschranke die
Auswertungshorizonte ins Unendliche -- und ein nicht-konstantes
Polynom kann denselben Wert nicht an beliebig vielen Stellen
annehmen. Das zwingt die Kettenglieder, konstant zu sein. Konstanten
aber sterben am nächsten Multiplikator, weil auch dort die
Ziffernschranke den Horizont linear anwachsen lässt, während die
Distributivität lineares Wertwachstum verlangt. Beide Mischformen
zerbrechen an derselben Stelle.

Der Satz hat inzwischen sogar eine dritte Dimension: Auch die
Bruch-Identität aus Kapitel~\ref{ch:erweiterungen} lässt sich weder
mit der Wert- noch mit der Gestaltseite verweben, und eine eigene
Arithmetik trägt sie nicht (ihre Addition kann die Eins
überschreiten). Die Seitentreue ist eine Trilogie: drei Identitäten,
zwei Rechenwelten, keine Mischung -- Beweise in
Anhang~\ref{ch:anhang-beweise}.

Die Pointe: Die Quelle der Starrheit ist die Ziffernbedingung
$0\le a_i\le i$ -- die allererste, unschuldigste Regel dieses Buchs.
Die Regel, mit der alles begann, ist genau die Regel, die verhindert,
dass Wert und Gestalt jemals in einer Arithmetik verschmelzen. Starr
sind nicht die Horizontregeln; starr sind die Identitäten.

\begin{oberflaeche}
Eine Hori-Zahl hat drei Identitäten: was sie wert ist, wie sie
aussieht, und welcher Bruch sie ist. Dieses Kapitel hat gezeigt,
dass die ersten beiden vollwertige, spiegelbildliche Rechenwelten
tragen, dass ein und dieselbe Eins in beiden Welten als Türsteher
arbeitet -- und, als bewiesener Satz: dass niemand in zwei Welten
gleichzeitig rechnen kann, während die dritte gar keine eigene
trägt. Man muss sich entscheiden. Das ist der tiefste Satz, den die
Horimetrik zu bieten hat: Konsistenz hat einen Preis, und der Preis
ist eine Identität.
\end{oberflaeche}

\chapter{Neun Zählfolgen aus der Horimetrik}
\label{ch:zaehlungen}

% Register: D8, D9, S9, S19, S22, S41, V11. Quellen: Anhang A.4-A.6,
% C.6, F.1-F.5, H, oeis/ENTWUERFE.md.

Theorien hinterlassen Fingerabdrücke: Zahlenfolgen, die beim Zählen
ihrer Objekte entstehen. Die Horimetrik hat im Lauf dieses Projekts
neun solcher Folgen erzeugt, und keine davon findet sich in der
OEIS, der großen Enzyklopädie der Zahlenfolgen -- was kein Beweis
für Neuheit ist, aber ein Indiz, dass hier bisher niemand gezählt
hat. Alle neun werden bei der OEIS eingereicht; die zugeteilten
A-Nummern tragen wir nach Vergabe hier nach.\footnote{Platzhalter:
A??????. Die Einreichungsentwürfe mit jargonfreien Definitionen und
Programmen liegen im Projektordner \texttt{oeis/}.}

\section{Die Steckbriefe}

\paragraph{Die Sprungstellen des Tiefengesetzes.}
$3,\ 15,\ 84,\ 533,\ 3883,\ 31998,\ 294875,\ 3006206,\ \ldots$
Die Schale, in der die Kommutator-Tiefe
(Satz~\ref{satz:tiefengesetz}) erstmals $-g$ erreicht.
Die Königin unserer Folgen: Für sie ist ein asymptotisches Gesetz
\emph{bewiesen} ($m_g/(g+2)!\to1$), und sie hat drei Vorhersagetests
bestanden -- den letzten exakt: $m_7=294875$ wurde prognostiziert,
bevor es berechnet war.

\paragraph{Die maximale Produkthöhe von Quadraten.}
$1,\ 6,\ 22,\ 87,\ 407,\ 2473,\ 19527,\ \ldots$
Wie groß ein Horizont werden kann, wenn man zwei randvolle
Gestalten der Kapazität $r$ multipliziert (Register D8). Inzwischen bewiesen: Die Folge wächst wie
$r!\cdot e^{2\sqrt r+O(\log r)}$ -- eine Fakultät mal einem sanft
explodierenden Faktor (Anhang~\ref{ch:anhang-beweise}); nur die
Feinstruktur des Logarithmus-Terms ist noch offen.

\paragraph{Die Schalenspringer.}
$1,\ 8,\ 60,\ 360,\ 2520,\ 21168,\ 211680,\ 2268000,\ \ldots$
Die Anzahl der Zahlen, die eine einzige führende Null in die
nächste Fakultätsschale hebt. Diese Folge hat uns das
$n!/2$-Debakel beschert (Kapitel~\ref{ch:sackgassen}): Drei Terme
lang sah sie aus wie eine Halbfakultät, beim vierten nicht mehr.
Ihr Komplement, die \emph{Nicht-Springer} $n!-D(n)$
($5,\ 16,\ 60,\ 360,\ 2520,\ 19152,\ \ldots$), zählt als eigene
Folge der Sammlung -- ausgezeichnet durch eine vollständig
bewiesene geschlossene Formel, die Ziffernformel aus
Anhang~\ref{ch:anhang-beweise}.

\paragraph{Die beidseitigen Fixpunkte.}
$4,\ 17,\ 88,\ 540,\ 3960,\ 32760,\ \ldots$
Zahlen, deren Darstellung unter beiden Aufräum-Operationen
gleichzeitig stabil ist (Abschnitt~\ref{sec:kommutator}); zählbar
als $n\cdot n!$ minus Schalenspringer, denn von den $n\cdot n!$
wertkompakten Elementen einer Schale scheitern an der
Strukturkompaktheit genau die ohne Anschlag-Ziffer -- und das sind
die Springer-Folgen.

\paragraph{Die Ein-Ziffern-Fakultäten.}
$1,\ 2,\ 5,\ 7,\ 11,\ 23,\ 29,\ 39,\ 59,\ 119,\ 143,\ \ldots$
Die $n$, für die $n!$ in seiner eigenen Schale eine einzelne Ziffer
ist -- vollständig charakterisiert als $n=(i+1)!/d-1$ mit
$1\le d\le i$, genau $i$ Stück pro Position
(Anhang~\ref{ch:anhang-beweise}); die Position $i=1$ liefert den
trivialen Anfangsterm $n=1$. Die einzige unserer Folgen, die
schon bei ihrer Entdeckung eine geschlossene Beschreibung mitbrachte
-- und der Lieferant von neun exakt eingetroffenen Vorhersagen.

\paragraph{Die Kompositions-Explosion.}
$1,\ 4,\ 23,\ 6541,\ 477\,149\,751,\ 26\,594\,371\,862\,819\,905,\ \ldots$
Wie hoch die Kapazität klettern kann, wenn man eine randvolle
Gestalt der Kapazität $r$ in eine ebensolche \emph{einsetzt}
(Kapitel~\ref{ch:strukturpolynome}). Die jüngste und mit weitem
Abstand am schnellsten wachsende Folge der Sammlung: Einsetzen
multipliziert die Grade und lässt die Koeffizienten explodieren.
Dank Extremalprinzip trotzdem mühelos berechenbar -- ein einziges
kanonisches Einsetzen pro Term. Ihre Asymptotik ist offen.

\paragraph{Die Interchange-Maxima und -Nullzähler.}
$0,\ 1,\ 6,\ 43,\ 351,\ 2873,\ 26443,\ 270291,\ \ldots$ und
$2,\ 5,\ 13,\ 23,\ 61,\ 111,\ 274,\ 564,\ \ldots$
Die beiden Wachstumsoperationen -- erst Null voran, dann werttreu
umziehen, oder umgekehrt -- vertauschen nicht. Den Unterschied
misst der \emph{Interchange-Defekt} $\Omega$: der Wert nach „erst
wachsen, dann liften“ minus der Wert nach „erst liften, dann
wachsen“. Dieser Defekt ist \emph{niemals negativ}: Erst wachsen,
dann liften liefert nie weniger. Das kleinste
lehrreiche Beispiel ist die Drei als $10$ in $H_2$: Erst wachsen:
$010$ in $H_3$ ist die \emph{Vier}; werttreu geliftet bleibt es die
Vier. Erst liften: die Drei wird zu $003$; eine Null davor ist
gratis, es bleibt die Drei. Also $\Omega=4-3=1$. Das war lange Vermutung und ist inzwischen
Satz -- der Beweis läuft über ein Multiplikationslemma mit scharfer
Regime-Grenze. Auch die Gleichheitsfälle sind inzwischen
charakterisiert: gleichgültig gegenüber der Reihenfolge ist eine
Zahl genau dann, wenn beide Beweisschritte scharf sind --
übertragsfrei plus scharfes Lemma -- und der Nullzähler gehorcht
einer daraus abgeleiteten semi-geschlossenen Zählformel (exakt
bestätigt für alle gemessenen Terme).

\section{Warum Folgen einreichen?}

Aus zwei Gründen, einem bescheidenen und einem strategischen. Der
bescheidene: Es ist ein kleiner, echter Dienst an der Mathematik --
saubere Definitionen, Programme, ein bewiesenes Gesetz. Der
strategische: Die OEIS ist der beste Neuheitsdetektor der Welt.
Wer irgendwo, in welchem Kontext auch immer, künftig
$3, 15, 84, 533$ ausrechnet, findet diese Einträge -- und damit
dieses Buch. Umgekehrt verlinken die Einträge hierher. Sollte eine
der Folgen doch eines Tages als bekannte Größe in Verkleidung
enttarnt werden, ist auch das ein Ergebnis: Dann wissen wir endlich
präzise, welcher Teil der Horimetrik alte Mathematik in neuen
Kleidern ist -- und welcher nicht.

\begin{oberflaeche}
Neun Zählfragen, neun Antwortreihen, die offenbar noch niemand
aufgeschrieben hat. Vielleicht, weil sie neu sind. Vielleicht nur,
weil nie jemand gefragt hat. Der Unterschied zwischen diesen beiden
Möglichkeiten ist genau das, was eine Einreichung klärt -- und
darum reichen wir ein.
\end{oberflaeche}

\chapter{Sackgassen: was nicht funktioniert hat und warum}
\label{ch:sackgassen}

% Quellen: verworfen.md X1-X9, counterexamples.md, lab-notes.
% Ton: erzählend, ehrlich, pro Sackgasse eine klare Lehre.

Forschungsberichte erzählen gern nur die Treffer. Dieses Kapitel
macht das Gegenteil: Es sammelt die Wege, die wir gegangen und
wieder zurückgegangen sind. Nicht aus Bußfertigkeit, sondern weil
die Sackgassen dieses Projekts mehr über die Horimetrik verraten
haben als mancher Satz -- und weil mindestens eine von ihnen sich
später als Umweg statt Irrweg entpuppte.

\section{Das $n!/2$-Debakel, oder: drei Datenpunkte sind kein Muster}

\begin{sackgasse}
\emph{Idee:} Die Schalenspringer-Folge zeigte bei $n=5,6,7$ die
Werte $60, 360, 2520$ -- exakt $n!/2$. Eine Formel!
\emph{Hindernis:} Der vierte Datenpunkt: $21168\ne20160$.
\emph{Lehre:} Muster aus wenigen Termen sind Hypothesen, keine
Ergebnisse.
\end{sackgasse}

Diese Geschichte hat ein bemerkenswertes Nachspiel: Viel später
fiel eine exakte Ziffernformel für dieselbe Folge vom Himmel -- und
sie erklärt rückwirkend, \emph{warum} die falsche Formel drei Terme
lang stimmte. Für alle drei Treffer reduziert sich die
Ziffernformel auf einen einzigen Summanden, und der beträgt jeweils
genau $n!/2$: bei $120$ und $5040$, weil sie in ihrer eigenen
Schale Ein-Ziffern-Darstellungen haben, bei $720$, weil die erste
Maximalziffer die Formel schon nach dem ersten beitragenden
Summanden abschneidet. Die Koinzidenz war kein Rauschen,
sondern ein echtes Strukturmerkmal -- nur eben eines der Ziffern
von $n!$, nicht der Folge selbst. So rehabilitieren sich Sackgassen
manchmal: Der Fehler von damals wurde zur Fußnote eines Satzes.

\section{Die Kleine-Welt-Falle, sechsmal}

Dieselbe Falle hat uns im Lauf des Projekts sechs Mal erwischt, in
immer neuen Verkleidungen: Eine Schranke des Kompaktifizierungs-Kommutators
sah bis zur siebten Welt wie $\pm1$ aus (sie ist in Wahrheit
beidseitig unbeschränkt); ihre Nachfolgerin, eine einseitige
Schranke, hielt Zufallsstichproben bis zur einundzwanzigsten Welt
stand und fiel erst durch \emph{konstruierte} Zeugen; ein Element
des Kompaktifizierungsmonoids wirkte bis zur vierten Welt
idempotent, und der kleinste Gegenzeuge wohnt in der fünften; und
das $n!/2$-Muster kennen Sie schon. Die Fälle fünf und sechs
ereigneten sich bei der oberen Seite des Tiefengesetzes an einem
einzigen Abend: Erst starb eine hübsche Dreieckszahlen-Hypothese am
eigenen Kandidaten-Sweep, noch ehe sie aufgeschrieben war -- dann
entpuppte sich die Schwelle, die dieser Sweep lieferte, ihrerseits
als Artefakt einer zu flachen Suche. Die gemeinsame Lehre ist
präziser als das übliche „mehr testen“:

\begin{center}
\emph{Im Horiraum wachsen die interessanten Gegenbeispiele
fakultätisch langsam.\\ Wer Schranken behauptet, muss konstruieren,
nicht stichproben.}
\end{center}

Übrigens spukt durch mehrere dieser Geschichten dieselbe Zahl:
die Sechs -- Ziel des kleinsten Schalensprungs (die Fünf, um eine
Null verlängert, wird zur Sechs), wertkleinster Trennzeuge des
Monoids, Anfang der ersten interessanten Schale. Sie ist der
Hausgeist dieses Buchs.

\section{Verworfene Anwendungsträume}

\begin{sackgasse}
\emph{Idee:} Hori-Adressen statt IPv6, Hori-Heaps statt Binärheaps.
\emph{Hindernis:} Der Horiraum ist innerhalb eines Horizonts
strukturell starr -- Adressräume brauchen lokale Elastizität, und
der Heap zahlt für die wachsende Verzweigung mit teuren
Kindervergleichen.
\emph{Lehre:} Die Stärke des Systems ist die globale
Fakultätsordnung, nicht lokale Flexibilität. Anwendungen müssen zur
Struktur passen, nicht umgekehrt.
\end{sackgasse}

\begin{sackgasse}
\emph{Idee:} Mächtigkeits-Argumente -- ist der Horiraum „größer“
als die natürlichen Zahlen?
\emph{Hindernis:} Er ist abzählbar; jede Nummerierung plättet ihn.
\emph{Lehre:} Der Mehrwert liegt nicht in der Kardinalität, sondern
in Struktur pro Wert. Ein Schachbrett wird nicht interessanter,
wenn man seine Felder durchnummeriert -- und nicht langweiliger.
\end{sackgasse}

Der dritte Anwendungstraum kam von außen, war der ehrgeizigste --
und starb am ersten Satz dieses Buchs.

\begin{sackgasse}
\emph{Idee:} Kryptographie aus Kapazität (externer Vorschlag,
Juli 2026). Eine geheime Gestalt wird durch $k$ führende Nullen
geschoben und wertkompakt ausgewertet; veröffentlicht wird nur der
nackte Wert. Der Angreifer, so die Hoffnung, scheitert am
Differenzüberlauf -- der Rückweg über Anti-Differenzen sprengt die
Kapazität -- und der Seitentreue-Satz versperrt jede hybride
Abkürzung.
\emph{Hindernis:} Der Rückweg braucht weder Anti-Differenzen noch
gemischte Arithmetik. Aus $y=P(N)$ holt fortgesetzte Division mit
Rest die Ziffern einzeln herunter: $c_0=y\bmod N$ ist \emph{exakt}
(denn $c_0\le r<N$), dann liefert $y_1=(y-c_0)/N$ die nächste Runde
mit Divisor $N{-}1$, und so fort -- eine Handvoll Divisionen, und
die „geheime“ Gestalt liegt offen. Das ist kein raffinierter
Angriff, sondern die Bijektion aus Kapitel~\ref{ch:horizonte} am
verschobenen Auswertungspunkt: Die Gewichte $1,N,N(N{-}1),\ldots$
bilden wieder ein Mischbasensystem, und die Ziffernschranken
bleiben brav unter den Divisoren.
\emph{Lehre:} Die Fakultätsordnung verwirbelt nicht, sie
\emph{sortiert} -- dieselbe Bijektivität, die das System so sauber
macht, verbietet ihm die kryptographische Härte. Zwei
Präzisierungen fürs Protokoll: Der Seitentreue-Satz verbietet
gemischte \emph{Halbringe}, nicht gemischte \emph{Algorithmen} --
ein Angreifer darf dividieren, tabellieren und Repräsentationen
wechseln, ohne je eine Arithmetik zu mischen. Und eine
Einwegfunktion braucht ein Geheimnis; die Horimetrik hat keins, sie
ist mit Absicht überall gleich hell. Was bleibt, ist die
bescheidenere Rolle: Das Horizontkalkül aus
Kapitel~\ref{ch:strukturpolynome} taugt als \emph{Messinstrument}
für Kapazitäts- und Rauschwachstum in fremden Konstruktionen --
nicht als eigene Härtequelle.
\end{sackgasse}

\section{Die tropische Fata Morgana -- und ihre Auflösung}

Die lehrreichste Sackgasse verdient den Schluss. Früh fiel auf,
dass die Horizonte sich unter Addition wie ein Maximum und unter
Multiplikation wie eine Summe verhalten \emph{wollen} -- das
Verhalten der sogenannten tropischen Algebra. Der direkte Versuch,
so zu rechnen, scheiterte an einem Überlaufproblem: Die Wertsumme
kann den Maximum-Horizont sprengen. Verworfen.

Monate später -- in der Zeitrechnung dieses Projekts: Stunden --
kehrte die Idee durch die Hintertür zurück. Der Horizont zerfällt
in einen Pflichtanteil und einen frei wählbaren Überschuss, und
\emph{auf dem Überschuss} funktioniert die tropische Arithmetik
tadellos: Der Pflichtanteil fängt alle Überläufe, das Gedächtnis
rechnet mit Maximum und Summe. Die Idee war nie falsch. Sie stand
nur im falschen Koordinatensystem.

\begin{oberflaeche}
Wenn dieses Kapitel eine Botschaft hat, dann diese: Verwerfen ist
keine Niederlage, sondern Buchführung. Von unseren neun
dokumentierten Sackgassen (das Arbeitsregister führt sie als X1
bis X9) wurde eine später zum Satz -- die tropische --, eine zur
Fußnote eines Satzes -- das $n!/2$-Debakel --, und die übrigen
wurden zu präzisen Warnschildern.
Nichts davon wäre passiert, hätten wir die Fehlversuche gelöscht
statt datiert.
\end{oberflaeche}

\chapter{Offene Fragen und ein Forschungsprogramm}
\label{ch:offen}

% Quelle: conjectures.md, Stand Version 0.7.

Bevor das letzte Kapitel die Theorie auf der Landkarte der
Mathematik verortet, ein ehrlicher Zwischenhalt: nicht, was
bewiesen wurde -- das haben die vorigen Kapitel getan --, sondern
was offen ist, und in welcher Reihenfolge es sich zu versuchen
lohnt.

\section{Die nächstliegenden Fragen}

\paragraph{Die Ziffernbilder der Fakultäten.} Die Ziffernformel für
die Schalenspringer hat die Frage verschoben: Nicht mehr „wie viele
Springer gibt es?“, sondern „wie schreibt sich $n!$ in seiner
eigenen Schale?“. Erste Muster liegen vor: Die Zahl der führenden
Nullen wächst wie die Umkehrfunktion der Fakultät, und die
Ein-Ziffern-Fälle sind inzwischen \emph{vollständig}
charakterisiert: Es sind exakt die $n=(i+1)!/d-1$ mit $1\le d\le i$,
genau $i$ Stück pro Position, und die Vorhersagen $n=29,39,59,119$
sowie $143,179,239,359,719$ trafen alle ein
(Anhang~\ref{ch:anhang-beweise}). Eine vollständige Theorie der
übrigen Ziffernbilder (Verteilung der Nichtnull-Anzahl, Position der
ersten Maximalziffer) fehlt weiterhin.

\paragraph{Die Feinstruktur der Interchange-Gleichheit.} Die
Positivität des Interchange-Defekts und die Charakterisierung der
Gleichheitsfälle sind bewiesen (Anhang~\ref{ch:anhang-beweise});
die Zählfolge $2$, $5$, $13$, $23$, $61$, $111$, $274$,
$564,\ldots$ gehorcht einer semi-geschlossenen Formel über die
scharfen Quotienten $q'$ aus dem Positivitätsbeweis, und
deren vollständige Ziffernentfaltung ist inzwischen bewiesen: pro
Ebene zwei Ungleichungen, ein wachsender Exzess, ein Abstieg
(Anhang~\ref{ch:anhang-beweise}). Die Asymptotik des
Nullzähler-Anteils ist inzwischen nach oben geklärt: Gleichgültigkeit gegenüber der Wachstumsreihenfolge ist
asymptotisch extrem selten -- der Anteil fällt mindestens wie
$e^{-(1+o(1))\,n\ln n}$, weil die Scharfheit die Ziffern ebenenweise
abwürgt. Offen bleibt die passende untere Schranke.

\paragraph{Die obere Seite des Tiefengesetzes.} Die untere Seite
ist inzwischen vollständig bewiesen -- Sprungstellen, Asymptotik
und, als letztes Stück während der Schlussredaktion, die exakte
Gleichheit: Die maximale Tiefe in Schale $m$ ist \emph{genau} die
Schalenlücke, und ein stets tauglicher Zeuge ist die
Terminaldarstellung des Maximalpolynom-Wertes (für Schale 3: die
Hausgeist-Zahl $(6,6)$).
Offen ist nur noch die obere Seite: Wie schnell wächst das
\emph{Maximum} des Kommutators? Hier ist Vorsicht geboten, und zwar
methodisch begründete: Exhaustiv gesichert ist nur der erste Sprung
($+2$ ab der neunten Welt). Alle späteren Schwellen sind bloß nach
oben beschränkt ($+3$ spätestens ab der vierzehnten, $+4$ spätestens
ab der zwanzigsten, $+5$ spätestens ab der sechsundzwanzigsten) --
jede Vertiefung der Zeugensuche hat sie bisher nach vorn verschoben,
und zwei voreilige Formelhypothesen sind an genau dieser Front
gestorben, eine davon noch vor ihrer Niederschrift.

\section{Die strukturellen Fragen}

\paragraph{Gekoppelte Arithmetiken.} Der Seitentreue-Satz verbietet
das Mischen der Identitäten; \emph{innerhalb} einer Seite ist die
Landschaft aber nicht klassifiziert. Die Produktkonstruktionen über
den Exzess liefern eine ganze Familie; der konservativ-kompakte
Halbring ist nachweislich keine davon. Gibt es weitere gekoppelte
Strukturen? Gibt es eine vollständige Klassifikation pro Seite?

\paragraph{Nachtrag: der nichtkommutative Mischfall ist geschlossen.}
Auch diese Lücke fiel während der Schlussredaktion: Verlangt man,
wie in jedem Halbring, \emph{beide} Distributivgesetze, so spiegelt
die Rechts-Distributivität die Kettenkonstruktion und ersetzt die
Kommutativität vollständig (Beweis in
Anhang~\ref{ch:anhang-beweise}). Was übrig bleibt, ist nur die
Kuriosität eines einseitig distributiven Mischsystems -- außerhalb
jeder Halbring-Axiomatik, und ehrlicherweise eher eine Fußnote als
eine Frage.

\paragraph{Nachtrag: die dreiseitige Seitentreue ist bewiesen.}
Während der Schlussredaktion fiel auch der letzte Fall
(Gestalt-Addition mit Bruch-Multiplikation): Die Kettengleichung
erzwingt die polynomiale Identität $2A(Y)=Y^{\underline{k}}$, die
den Koeffizienten $\tfrac12$ verlangt -- Ziffern sind aber ganz.
Zusammen mit den zwei zuvor entschiedenen Fällen gilt: Es gibt genau
zwei Rechenwelten (Wert und Gestalt), und der Bruch trägt keine
dritte und mischt sich in keine. Details in
Anhang~\ref{ch:anhang-beweise}. Offen bleibt hier also -- nichts
mehr; die Frage wandert als bewiesener Satz in
Kapitel~\ref{ch:arithmetik}.

\paragraph{Formalisierung.} Der Seitentreue-Satz ist elementar
genug für eine vollständige Formalisierung in Lean; die fallenden
Fakultäten liegen in Mathlib bereits vor. Das wäre der Goldstandard
der Absicherung -- und ein Beitrag zur Formal-Methods-Gemeinde
gleich mit.

\section{Woran sich die Theorie messen lassen will}

Zum Schluss die Frage hinter allen Fragen: Wann wäre die Horimetrik
\emph{gescheitert}, wann \emph{gelungen}? Gescheitert, wenn sich
ein bestehender Formalismus findet, in dem Horizont, Dualität und
Seitentreue nur Umbenennungen sind -- die Literatursuche hat ihn
nicht gefunden, aber Suchen beweisen nichts. Gelungen, wenn eines
von dreien eintritt: wenn die Zählfolgen dieses Buchs anderswo
unabhängig auftauchen; wenn der Seitentreue-Satz oder das
Tiefengesetz in fremder Sprache neu bewiesen wird; oder wenn
jemand, der dieses Buch nie gelesen hat, an den Ziffern von $n!$
in der eigenen Schale dieselbe Freude findet wie wir.

\begin{oberflaeche}
Eine Theorie ist fertig, wenn sie keine Fragen mehr erzeugt. Nach
diesem Maßstab ist die Horimetrik erfreulich unfertig -- und genau
so soll dieser Zwischenhalt enden: nicht mit einem Punkt, sondern
mit einem Doppelpunkt. Was bleibt, ist die Frage nach der Adresse:
Wo auf der Landkarte wohnt das alles?
\end{oberflaeche}

\chapter{Einordnung: Wo die Horimetrik wohnt}
\label{ch:einordnung}

% Register: S44, V13. Quellen: Anhang J, experiments_einordnung.py.
% Schlusskapitel: erst die Fronten (K8), dann die Landkarte.

Ein Buch über eine neue Theorie schuldet seinen Lesern am Ende eine
Adresse: Wo auf der Landkarte der Mathematik liegt das alles? Die
Antwort hat sich erst ganz zum Schluss der Arbeit vollständig
zusammengesetzt, und sie ist besser, als wir zu hoffen wagten. Es
gibt nämlich drei mögliche Ausgänge für ein Projekt wie dieses. Ein
\emph{unbekannter} Grundraum wäre verdächtig -- vermutlich hätte man
nur schlecht gesucht. Eine \emph{bekannte vollständige Theorie} wäre
das Ende -- man hätte alte Mathematik neu angemalt. Der beste Fall
ist der dritte: ein bekannter Grundraum, auf dem neue Operationen
und neue Sätze stehen. Genau dort sind wir gelandet.

\section{Die Bewohner sind alte Bekannte}

Ergänzt man die Ziffern einer Hori-Zahl um eine führende Null,
$e=(0,a_1,\ldots,a_n)$, so gilt an jeder Stelle $0\le e_j<j$ -- und
das ist Wort für Wort die Definition einer
\emph{Inversionsfolge} der Länge $n{+}1$. Inversionsfolgen stehen
in klassischer Bijektion zu den Permutationen: $e_j$ zählt, wie
viele frühere Einträge einer Permutation größer sind als der
$j$-te. Also:

\begin{satz}[Register S44]
$H_n$ steht in Bijektion zu den Inversionsfolgen der Länge $n{+}1$
und damit zu den Permutationen von $n{+}1$ Elementen. Unter dieser
Brille wird der Strukturhorizont zur Permutationsstatistik
\[
\sigma(P_a)=n-\min_{e_j>0}\bigl((j{-}1)-e_j\bigr),
\]
und seine Verteilung ist exakt bekannt: Auf jedem Niveau
$1\le r\le n$ liegen genau $r\cdot r!$ Elemente (die Formel gilt
für $a\ne0$; die Null hat $\sigma=0$).
\end{satz}

Die Fakultätsgröße der Horizonte war also nie ein Zufall: Die
Bewohner der Hori-Welten \emph{sind} Permutationscodes. Das nimmt
der Theorie nichts -- es sagt präzise, was sie ist: Die Horimetrik
untersucht, wie diese wohlbekannten Codes ihren \emph{Wert}, ihre
\emph{Gestalt} und ihre \emph{Identität} verändern, wenn ihr Raum
wächst. Umzüge, Projektoren, Seitentreue, Kapazität -- all das ist
auf dem alten Grundraum neu.

Die Statistik selbst hat eine hübsche Lesart: $(j{-}1)-e_j$ ist der
Abstand einer Ziffer von ihrer oberen Schranke. Die Literatur kennt
die Extremfälle als \emph{tight entries} (Ziffern, die ganz oben
anliegen); der Strukturhorizont verfeinert das zu einem Maß: Wie
nahe kommt \emph{irgendeine} Ziffer ihrem Anschlag? Ob diese
Statistik unter der Bijektion einer klassischen
Permutationsstatistik entspricht, ist offen (Register V13) -- ihre
Verteilung $r\cdot r!$ ist jedenfalls fast verdächtig sauber.

\section{Die Ordnung: ein Gitter, das es schon gibt}

Vergleicht man Ziffernfolgen stellenweise, wird jeder Horizont zu
einem Produkt endlicher Ketten -- einem \emph{distributiven
Gitter}. In dieser Ordnung ist die randvolle Gestalt $M_r$ ein
gewöhnliches Element, und die Kapazitätsklasse
$\mathcal F_r=\{P:\sigma(P)\le r\}$ ist schlicht das
\emph{Hauptideal} darunter: alles, was unter $M_r$ liegt. Das ist
die sauberste abstrakte Beschreibung des Strukturhorizonts, die wir
kennen: $\sigma(P)$ ist die erste Stufe der Kette
$M_0\le M_1\le M_2\le\cdots$ (mit $M_0=0$), deren Hauptideal $P$
enthält -- und
das Extremalprinzip aus Kapitel~\ref{ch:strukturpolynome} wird in
dieser Sprache beinahe eine Tautologie. Dieselbe stellenweise
Ordnung auf Inversionsfolgen wurde jüngst als \emph{middle order}
auf Permutationen untersucht, zwischen schwacher und
Bruhat-Ordnung. Ob unsere Projektoren $K_V$, $K_S$ in diesem Gitter
monoton sind, ob der Kommutator ein Ordnungsdefekt ist -- das sind
ausformulierte Anschlussfragen an ein existierendes Gebiet.

\section{Der Kalkül: alte Operatoren, neue Buchhaltung}

Die fallenden Fakultäten sind die Hausbasis der
\emph{finiten Operatorrechnung}: $\Delta$ spielt dort die Rolle der
Ableitung, und die umbrale Algebra studiert seit Jahrzehnten
Delta-Operatoren und ihre Basisfolgen. Strukturpolynome und
Horizontableitung sind also bestens verortet. Was die Horimetrik
hinzufügt, ist eine Frage, die dort nicht im Mittelpunkt steht:
nicht \emph{wie} ein Operator ein Polynom verändert, sondern
\emph{wie viel minimale Strukturkapazität} das kostet. Das
Horizontkalkül -- $\Gamma_+$, $\Gamma_\times$, $\Gamma_\Delta$,
$\Gamma_E$, $\Gamma_\circ$, alle exakt und alle am selben
kanonischen Objekt angenommen -- ist eine Kapazitätsbuchhaltung
über einem etablierten Operatorgerüst. Zusammen mit der
Zählbreiten-Deutung aus Kapitel~\ref{ch:strukturpolynome} ist das
der Kandidat für einen eigenständigen Fachaufsatz ganz ohne
Hori-Erzählung.

\section{Das Forschungsprogramm dahinter}

Drei weitere Nachbarschaften haben wir gesichtet, aber nicht mehr
betreten; sie seien als Programm notiert. \emph{Erstens:} Ersetzt
man die Ziffernschranken $0\le a_i\le i$ durch beliebige
Kapazitätsprofile $0\le a_i<s_i$, entstehen die in der Kombinatorik
untersuchten $s$-Inversionsfolgen samt ihrer Geometrie
(Lecture-Hall-Polytope). Eine \emph{$s$-Horimetrik} wäre eine
zweite Theorie -- welche Profile eine Polynombasis, eine
Wert-Gestalt-Dualität, eine Seitentreue erlauben, ist vollständig
offen. \emph{Zweitens:} Die Rook-Theorie schreibt ihre Polynome in
genau unserer Basis; welche Gestalten Rook-Polynome von Brettern
sind und ob $M_r$ ein natürliches Brett besitzt, wäre ein
kombinatorischer Zugang zu unseren Produktformeln.
\emph{Drittens:} Die Zählsemantik der Gestalten -- gefärbte
geordnete Injektionen -- ist der Rohstoff der kombinatorischen
Spezies; dort könnte eine gemeinsame Sprache für alle
Interpretationen dieses Buchs liegen.

\begin{oberflaeche}
Das Fazit des Buchs, in drei Zeilen: Die Bewohner der Horimetrik
sind bekannte Objekte -- Permutationscodes in einem bekannten
Gitter, beschrieben in der Hausbasis der Differenzenrechnung. Die
Operationen darauf -- zwei Arten zu wachsen, zwei Arten
aufzuräumen, eine Kapazitätsbuchhaltung, eine Zählbreite -- sind
neu, und die Sätze über sie auch. Ein bekannter Grundraum mit
neuen Operationen und neuen Sätzen: Für eine kleine Theorie, die
mit dem Umdrehen einer Ziffernordnung im Dunkeln begann, ist das
vermutlich der beste Platz, den die Landkarte zu vergeben hat.
\end{oberflaeche}

\chapter*{Literatur und Einordnung}
\addcontentsline{toc}{chapter}{Literatur und Einordnung}
\markboth{LITERATUR UND EINORDNUNG}{LITERATUR UND EINORDNUNG}

Dieses Verzeichnis ist bewusst kommentiert. Die Bausteine der
Horimetrik -- Mischbasensysteme, fallende Fakultäten,
Inversionsfolgen, das Differenzenkalkül -- sind etablierte
Mathematik, und dieses Buch behauptet nichts anderes. Jeder Eintrag
nennt deshalb zweierlei: welchen Baustein er beisteuert, und was
dort jeweils \emph{nicht} steht -- also die Stelle, an der die
Horimetrik über ihre Quelle hinausgeht. Wer die Einordnung aus
Kapitel~\ref{ch:einordnung} nachprüfen will, findet hier die
Belege.

% thebibliography setzt sonst ein eigenes \chapter* ("Literaturverzeichnis")
% auf eine neue Seite; wir haben oben schon eine Kapitelüberschrift.
\begingroup
\makeatletter
\renewcommand{\chapter}[2]{}%
\renewcommand{\@mkboth}[2]{}%
\makeatother
\begin{thebibliography}{10}

\bibitem{cantor} G.~Cantor: \emph{Über die einfachen Zahlensysteme},
Zeitschrift für Mathematik und Physik 14 (1869), 121--128.
--- Die Fakultätsbasis für Brüche; das ist exakt der bruchtreue
Grenzraum der dritten Treppe (Kapitel~\ref{ch:erweiterungen}),
nicht aber die gestalttreue Welt oder das Zusammenspiel.

\bibitem{knuth} D.~E. Knuth: \emph{The Art of Computer Programming},
Band~2, 3.~Aufl., Addison--Wesley 1997, Abschnitt 4.1.
--- Mischbasensysteme und das Fakultätszahlensystem (Factoradic);
dort wachsen die Basen zur signifikanten Seite, führende Nullen sind
wertneutral.

\bibitem{lehmer} D.~H. Lehmer: \emph{Teaching combinatorial tricks
to a computer}, Proc. Sympos. Appl. Math. 10 (1960), 179--193.
--- Lehmer-Codes und Inversionstabellen; erklären die
Fakultätsgröße der Horizonte und die Ziffernschranke $0\le a_i\le i$,
nicht aber die horizontabhängige Auswertung.

\bibitem{gkp} R.~L. Graham, D.~E. Knuth, O.~Patashnik:
\emph{Concrete Mathematics}, 2.~Aufl., Addison--Wesley 1994, Kap.~2.
--- Fallende Fakultäten und Differenzenkalkül; das Werkzeug der
Horizontableitung.

\bibitem{cahen} P.-J. Cahen, J.-L. Chabert: \emph{Integer-Valued
Polynomials}, AMS Surveys and Monographs 48, 1997.
--- Polynome mit ganzzahligen Werten in der Binomialbasis; die
Horimetrik nutzt nur den positiven Kegel und wertet am laufenden
Horizont aus.

\bibitem{rigo} M.~Rigo, M.~Stipulanti: \emph{Revisiting regular
sequences in light of rational base numeration systems},
arXiv:2103.16966.
--- Abstrakte Numerationssysteme, in denen führende Nullen über die
Radix-Ordnung wertsensitiv sind: der nächste bekannte Nachbar der
Nullen-Semantik, dort formalsprachlich statt algebraisch behandelt.

\bibitem{kudr} G.~Kudryavtseva, V.~Mazorchuk: \emph{On Kiselman's
semigroup}, arXiv:math/0511374; O.~Ganyushkin, V.~Mazorchuk:
\emph{On Kiselman quotients of 0-Hecke monoids}, arXiv:1006.0316.
--- Monoide aus idempotenten Erzeugern mit Absorptionsrelationen;
Nachbarschaft des sechselementigen Kompaktifizierungsmonoids.

\bibitem{simon} D.~Simon: \emph{Mixed radix numeration bases:
Horner's rule, Yang-Baxter equation and Furstenberg's conjecture},
arXiv:2405.19798.
--- Lokale Basiswechsel in Mischbasensystemen; Kandidatenwerkzeug
für den Interchange-Defekt.

\bibitem{invseq} S.~Corteel, M.~A.~Martinez, C.~D.~Savage,
M.~Weselcouch: \emph{Patterns in Inversion Sequences I},
arXiv:1510.05434.
--- Inversionsfolgen ($0\le e_j<j$) und ihre Bijektion zu
Permutationen; unser $H_n$ ist als Menge genau dieser Raum
(Kapitel~\ref{ch:einordnung}). Was dort nicht steht: die umgedrehte
Wertung, die Horizontübergänge und der Strukturhorizont als
Statistik.

\bibitem{lecturehall} C.~D. Savage, M.~J. Schuster: \emph{Ehrhart
series of lecture hall polytopes and Eulerian polynomials for
inversion sequences}, J. Combin. Theory Ser.~A 119 (2012),
850--870.
--- $s$-Inversionsfolgen mit beliebigen Kapazitätsprofilen und ihre
Geometrie; der Rahmen für das $s$-Horimetrik-Programm aus
Kapitel~\ref{ch:einordnung}.

\bibitem{tightentries} S.~Chern, S.~Fu, Z.~Lin: \emph{Burstein's
permutation conjecture, Hong and Li's inversion sequence
conjecture, and restricted Eulerian distributions},
arXiv:2209.12137.
--- Statistiken auf Inversionsfolgen, darunter die \emph{tight
entries} ($e_i=s_i-1$, Statistik $\mathrm{tig}$); der
Strukturhorizont verfeinert diese Extremfall-Zählung zu einem
Abstandsmaß (Kapitel~\ref{ch:einordnung}).

\bibitem{middleorder} M.~Bouvel, L.~Ferrari, B.~E.~Tenner:
\emph{Between weak and Bruhat: the middle order on permutations},
arXiv:2405.08943.
--- Die koordinatenweise Ordnung auf Inversionscodes als Gitter
zwischen schwacher und Bruhat-Ordnung; der Rahmen, in dem
$\mathcal F_r$ ein Hauptideal und $\sigma$ eine Idealstufe wird.

\bibitem{umbral} K.~Beauduin: \emph{Operational Umbral Calculus},
arXiv:2407.16348.
--- Moderne finite Operatorrechnung: Delta-Operatoren und ihre
Basisfolgen. Der Rahmen für $\Delta$ und die Strukturpolynome; was
dort nicht steht, ist die Kapazitätsbuchhaltung $\Gamma_T$.

\bibitem{ikenmeyerpak} C.~Ikenmeyer, I.~Pak: \emph{What is in \#P
and what is not?}, FOCS 2022, arXiv:2204.13149.
--- Charakterisiert die \emph{binomial-good}-Polynome
(nichtnegative ganzzahlige Binomialbasis-Koeffizienten) als die
\emph{relativierenden} polynomiellen \#P-Abschlussoperationen und
formuliert die Binomial-Basis-Vermutung (unrelativiert gilt
dieselbe Charakterisierung nur vermutungsweise), deren univariate
Fassung bereits P$\,\ne\,$NP implizieren würde. Die Gestalten der Horimetrik sind wegen
$\ff Xd=d!\binom Xd$ eine echte Teilklasse der binomial-good-Polynome
(Kapitel~\ref{ch:strukturpolynome}); was dort \emph{nicht} steht,
ist die quantitative Schicht: die Höhe $\max(d+c_d)$, die
$(r{+}1)!$-Hierarchien und deren Wachstumsgesetze.

\bibitem{oeis} OEIS Foundation Inc.: \emph{The On-Line Encyclopedia
of Integer Sequences}, \texttt{https://oeis.org}.
--- Referenzdatenbank für die neun Zählfolgen dieses Buchs;
A-Nummern werden nach Vergabe nachgetragen.

\end{thebibliography}
\endgroup


\appendix
\chapter{Referenz: Implementierung, Datentafeln, Registerstand}
\label{ch:anhang-referenz}

\section*{Die Referenzimplementierung}

Alle Enumerationen und Verifikationen dieses Buchs beruhen auf
einer kleinen, bewusst unoptimierten Python-Implementierung
(\texttt{horimetrik\_reference.py}) mit Selbsttests: Codierung und
Decodierung, Struktur\-polynome, beide Erweiterungen, beide
Kompaktifizierungen, Strukturarithmetik. Ergänzt wird sie durch
datierte Experimentskripte, deren Ergebnisse im Arbeitsregister mit
Status (bewiesen / vermutet / widerlegt) geführt werden. Der
Grundsatz des Projekts: Enumeration erzeugt Vermutungen und
Gegenbeispiele, niemals Sätze -- und jede behauptete Schranke muss
konstruierte Extremfälle überleben, nicht nur Stichproben.

\section*{Die neun Zählfolgen}

\begin{center}
\footnotesize
\begin{tabular}{lll}
\toprule
Folge & Anfangswerte & Status\\
\midrule
Sprungstellen $m_g$ ($g{\ge}1$) & $3,15,84,533,3883,31998,294875,\ldots$ &
Tiefengesetz, Asymptotik bewiesen\\
$\beta(r,r)$ ($r{\ge}1$) & $1,6,22,87,407,2473,19527,\ldots$ & Asymptotik bewiesen\\
Springer $D(n)$ ($n{\ge}3$) & $1,8,60,360,2520,21168,211680,2268000$ &
Ziffernformel bewiesen\\
Nicht-Springer $N(n)$ ($n{\ge}3$) & $5,16,60,360,2520,19152,151200,\ldots$ &
Ziffernformel bewiesen\\
Fixpunkte $F(n)$ ($n{\ge}2$) & $4,17,88,540,3960,32760$ & $=n\cdot n!-D(n)$ (bewiesen)\\
Ein-Ziffern-Fakultäten & $1,2,5,7,11,23,29,39,59,119,\ldots$ &
vollständig charakterisiert\\
$\max\Omega$ ($n{\ge}1$) & $0,1,6,43,351,2873,26443,270291$ & $\Omega\ge0$
bewiesen\\
$\#\{\Omega=0\}$ ($n{\ge}1$) & $2,5,13,23,61,111,274,564$ &
Gleichheitskriterium bewiesen\\
$\Gamma_\circ(r,r)$ ($r{\ge}1$) & $1,4,23,6541,477149751,\ldots$ &
Extremalprinzip bewiesen\\
\bottomrule
\end{tabular}
\end{center}

Keine dieser Folgen war zum Zeitpunkt der Niederschrift in der
OEIS; die Einreichung ist vorbereitet (A-Nummern werden hier
nachgetragen).

\section*{Bestätigte Vorhersagen}

\begin{enumerate}
\item Tiefe $-3$ des Kommutators erstmals nahe Schale $100$
      (exakt: 84 als Sprungstelle, bestätigt im Bereich bis 200).
\item $m_6=31998$ bei prognostizierten $\approx31930$.
\item $m_7=294875$ -- exakt wie vorhergesagt.
\item Ein-Ziffern-Fakultäten bei $n=29,39,59,119$ und
      $143,179,239,359,719$ aus der (inzwischen vollständigen)
      Charakterisierung $n=(i{+}1)!/d-1$ -- alle neun bestätigt.
\end{enumerate}

\section*{Registerstand bei Redaktionsschluss dieser Fassung}

44 bewiesene Sätze -- darunter die Grenzraum\-sätze der drei
Treppen, die zwei dualen Halbringe, der Dualitäts\-satz und der
Seitentreue-Satz, das Kompaktifizierungs\-monoid, das
Tiefengesetz, die Ziffernformel und die Interchange-Positivität
--, dazu 9 dokumentierte Sackgassen, 3 widerlegte eigene
Vermutungen, 9 unbekannte Zählfolgen und 4 bestätigte
Vorhersage-Gruppen (12 Einzeltreffer, Liste oben). Die jeweils aktuelle Fassung von Register,
Vermutungen, Gegenbeispielen und Changelog liegt den Arbeitsdateien
des Projekts bei.

\chapter{Die Beweise}
\label{ch:anhang-beweise}

Die Kapitel dieses Buchs erzählen die Beweisideen; hier stehen die
Beweise selbst, nach Kapiteln geordnet. Notation wie im Haupttext;
$E_n$ bezeichnet die Codierung $V_n^{-1}$, $\mu(x)$ den
Minimalhorizont, $\sigma(P)=\max_{c_d>0}(d+c_d)$ (mit
$\sigma(0)=0$) den Strukturhorizont.

\section*{Zu Kapitel 3: die Grenzraumsätze}

\begin{satz}[Satz~\ref{satz:grenzen}, vollständig]
$\varinjlim(H_n,Z_n)\cong\Struk$ und
$\varinjlim(H_n,L_n)\cong\Nn$.
\end{satz}

\begin{proof}
\emph{Strukturtreu.} Für $m\ge n$ sei $Z_{n,m}$ das Voranstellen von
$m-n$ Nullen; offenbar $Z_{m,k}\circ Z_{n,m}=Z_{n,k}$, das System
ist gerichtet. Im direkten Limes gelten $a\in H_n$ und $b\in H_m$
(o.B.d.A.\ $n\le m$) als gleich, wenn sie in einem gemeinsamen
späteren Horizont dasselbe Bild haben, also
$0^{k-n}a=0^{k-m}b$, was äquivalent zu $b=0^{m-n}a$ ist. Die
Klassen sind also Ziffernfolgen modulo führender Nullen. Die
Abbildung $[a]\mapsto P_a$ ist wohldefiniert (eine führende Null
fügt den Koeffizienten $0$ am höchsten Grad hinzu), injektiv (die
Ziffernfolge in $H_n$ ist die auf Länge $n$ aufgefüllte, von rechts
gelesene Koeffizientenfolge; gleiche Polynome erzwingen also
$b=0^{m-n}a$) und surjektiv (jedes $P\in\Struk$ ist ab
$H_{\sigma(P)}$ darstellbar).

\emph{Werttreu.} Sei $L_{n,m}=E_m\circ V_n$. Wegen
$V_m\circ E_m=\mathrm{id}$ gilt
$L_{m,k}\circ L_{n,m}=E_k\circ V_n=L_{n,k}$. Zwei Elemente werden
genau dann identifiziert, wenn $E_k(V_n(a))=E_k(V_m(b))$ für ein
$k$, und da $E_k$ injektiv ist, genau bei $V_n(a)=V_m(b)$: Die
Klassen sind die Fasern der Wertabbildung. $[a]\mapsto V_n(a)$ ist
wohldefiniert, injektiv und surjektiv (jedes $x$ tritt als
$V_{\mu(x)}(E_{\mu(x)}(x))$ auf).
\end{proof}

\section*{Zu Kapitel 3: die dritte Treppe}

\begin{satz}[Cantor-Koordinate und Rechts-Erweiterung]
Mit $F(a)=\sum_i a_i/(i+1)!$ gilt $V_n(a)=(n+1)!\cdot F(a)$. Das
Rechts-Anhängen $R(a_1,\ldots,a_n)=(a_1,\ldots,a_n,0)$ ist
bruchtreu, skaliert den Wert exakt mit $n+2$, und
$\varinjlim(H_n,R_n)$ ist die Menge der endlichen
Fakultätsbrüche in $[0,1)$.
\end{satz}

\begin{proof}
Termweise ist $a_i\,(n+1)!/(i+1)!=(n+1)!\cdot a_i/(i+1)!$, also die
Faktorisierung. Beim Rechts-Anhängen behalten alle Ziffern Position
und Nenner, die neue Null trägt nichts bei: $F$ invariant; der Wert
ist $(n+2)!\,F=(n+2)\cdot(n+1)!\,F$. Im Limes werden genau
Ziffernfolgen identifiziert, die sich um angehängte Nullen
unterscheiden; jede Klasse trägt genau einen endlichen Bruch, und
jeder endliche Fakultätsbruch tritt auf.
\end{proof}

\section*{Zu Kapitel 4: die Horizontableitung}

\begin{satz}[$\Delta$-Kalkül]
Die Wertänderung bei strukturtreuer Erweiterung ist
$\Delta P(n+1)$ mit $\Delta P(X)=P(X{+}1)-P(X)$ und
$\Delta\,\ff Xd=d\,\ff X{d-1}$; eine führende Null ist genau dann
wertneutral, wenn $P$ konstant ist.
\end{satz}

\begin{proof}
Unter $Z$ bleibt $P$ erhalten und der Auswertungspunkt wandert von
$n{+}1$ zu $n{+}2$; die Änderung ist per Definition
$\Delta P(n{+}1)$. Zur Differenzformel:
$\ff{(X{+}1)}d=(X{+}1)\,X\cdots(X{-}d{+}2)$ und
$\ff Xd=X\cdots(X{-}d{+}2)\,(X{-}d{+}1)$ haben den gemeinsamen
Faktor $X(X{-}1)\cdots(X{-}d{+}2)=\ff X{d-1}$; die verbleibenden
Faktoren sind $X{+}1$ bzw.\ $X{-}d{+}1$, ihre Differenz ist $d$,
also $\Delta\,\ff Xd=d\,\ff X{d-1}$.
Wertneutralität: $\Delta P=0$ für alle Auswertungspunkte im
zulässigen Bereich erzwingt, da Polynome mit unendlich vielen
Nullstellen verschwinden, $\Delta P\equiv0$, also $P$ konstant;
umgekehrt klar. Für den einzelnen Schritt: $\Delta P(n{+}1)=0$ mit
nichtnegativen Koeffizienten von $\Delta P$ (alle
Basisdifferenzen haben nichtnegative Koeffizienten) und
Auswertungspunkt im positiven Bereich -- wegen
$\deg\Delta P\le n-2$ sind alle $\ff{(n{+}1)}d$ dort strikt
positiv -- erzwingt ebenfalls $\Delta P=0$.
\end{proof}

\section*{Zu Kapitel 4: das Horizontkalkül}

Wir schreiben $P\preceq Q$, wenn jeder Koeffizient von $P$ in der
fallenden Fakultätsbasis höchstens der entsprechende Koeffizient von
$Q$ ist, und $M_r(X)=\sum_{d=0}^{r-1}(r{-}d)\,\ff Xd$ für die maximal
gefüllte Gestalt des Horizonts $r$.

\begin{satz}[Extremalprinzip; Register S39]
Sei $T$ ein $k$-stelliger Operator auf dem Strukturhalbring, der in
jedem Argument koeffizientenmonoton ist (aus $P\preceq P'$ folgt
$T(\ldots,P,\ldots)\preceq T(\ldots,P',\ldots)$). Dann gilt
\[
\Gamma_T(r_1,\ldots,r_k)
:=\max\{\sigma(T(P_1,\ldots,P_k)) \mid \sigma(P_i)\le r_i\}
=\sigma\bigl(T(M_{r_1},\ldots,M_{r_k})\bigr).
\]
Jeder aus Addition, Multiplikation und $\Delta$ zusammengesetzte
Operator ist koeffizientenmonoton.
\end{satz}

\begin{proof}
$\sigma(P)=\max_{c_d>0}(d+c_d)$ ist monoton bezüglich $\preceq$, und
$\sigma(P)\le r$ ist äquivalent zur koeffizientenweisen Schranke
$c_d\le r-d$, also zu $P\preceq M_r$. Aus der Monotonie von $T$
folgt $T(P_1,\ldots,P_k)\preceq T(M_{r_1},\ldots,M_{r_k})$ und damit
$\sigma(T(P_1,\ldots,P_k))\le\sigma(T(M_{r_1},\ldots,M_{r_k}))$;
wegen $\sigma(M_{r_i})=r_i$ wird die Schranke angenommen. Zur
Zusatzaussage: Addition und $\Delta$ sind koeffizientenweise linear
mit nichtnegativen Konstanten ($\Delta\,\ff Xd=d\,\ff X{d-1}$); für
die Multiplikation sind die Linearisierungskoeffizienten
$\binom pj\binom qj\,j!$ der Produktformel nichtnegativ, und die
Bilinearität überträgt die Monotonie. Kompositionen monotoner
Operatoren sind monoton.
\end{proof}

\begin{satz}[Differenzüberlauf; Register S40]
Für $r\ge2$ gilt
$\Gamma_\Delta(r)=\bigl\lfloor(r+1)^2/4\bigr\rfloor-1$,
angenommen bei der Ein-Ziffern-Gestalt $P=(r{-}d)\,\ff Xd$ mit
$d=\lfloor(r+1)/2\rfloor$.
\end{satz}

\begin{proof}
Nach dem Extremalprinzip ist $\Gamma_\Delta(r)=\sigma(\Delta M_r)$.
$\Delta M_r$ hat die Koeffizienten $b_k=(k{+}1)(r{-}k{-}1)$ für
$k=0,\ldots,r{-}2$, also
\[
\sigma(\Delta M_r)
=\max_{0\le k\le r-2}\bigl[k+(k{+}1)(r{-}k{-}1)\bigr]
=\max_{0\le k\le r-2}(k{+}1)(r{-}k)\;-\;1 .
\]
Mit $u=k{+}1\in\{1,\ldots,r{-}1\}$ ist $u(r{+}1{-}u)$ das Produkt
zweier Zahlen mit fester Summe $r{+}1$, maximal bei
$u=\lfloor(r{+}1)/2\rfloor$ mit Wert
$\lfloor(r{+}1)^2/4\rfloor$; der Bereich ist für $r\ge2$ nicht leer.
Die Ein-Ziffern-Gestalt $P=(r{-}d)\,\ff Xd$ hat $\sigma(P)=r$ und
$\sigma(\Delta P)=d-1+d(r{-}d)=d(r{+}1{-}d)-1$, maximal beim selben
$d$ mit demselben Wert. (Exhaustiv bestätigt für alle $(r{+}1)!$
Gestalten bis $r=8$: Werte $1,3,5,8,11,15,19$.)
\end{proof}

\begin{satz}[Kompositionsabschluss; Register S41]
Für $P,Q$ mit nichtnegativen ganzzahligen Koeffizienten in der
fallenden Fakultätsbasis hat auch $P\circ Q$ nichtnegative
ganzzahlige Koeffizienten. Die Komposition ist in beiden Argumenten
koeffizientenmonoton; nach dem Extremalprinzip gilt
$\Gamma_\circ(r,s)=\sigma(M_r\circ M_s)$.
\end{satz}

\begin{proof}
Wegen $P\circ Q=\sum_d c_d(P)\,(Q)^{\underline d}$ mit
$(Q)^{\underline d}=Q(Q{-}1)\cdots(Q{-}d{+}1)$ genügt es,
$(Q)^{\underline d}$ als zulässige Gestalt nachzuweisen. Schreibe
$Q=\sum_k a_k\ff Xk$ und wähle paarweise disjunkte Mengen $T_k$ mit
$|T_k|=a_k$ (\emph{Schablonen der Stelligkeit $k$}). Eine
\emph{$Q$-Struktur} auf $[x]$ sei ein Paar $(t,\varphi)$ aus einer
Schablone $t\in T_k$ und einer Injektion
$\varphi\colon[k]\hookrightarrow[x]$; ihre Anzahl ist
$\sum_k a_k\,\ff xk=Q(x)$. Also zählt $Q(x)^{\underline d}$ die
$d$-Tupel paarweise verschiedener $Q$-Strukturen auf $[x]$.

Klassifiziere die Tupel nach der Vereinigung $U\subseteq[x]$ der
Bilder aller beteiligten Injektionen. Die Anzahl $N_m$ der Tupel,
deren Bildvereinigung eine feste $m$-Menge ganz ausschöpft, hängt
nur von $m$ ab, und für alle $x\in\Nn$ gilt die Polynomidentität
$Q(x)^{\underline d}=\sum_m N_m\binom xm$. Die symmetrische Gruppe
$S_m$ wirkt auf den ausschöpfenden Tupeln durch Umbenennung der
Grundmenge, $\pi\cdot(t,\varphi)=(t,\pi\circ\varphi)$, und diese
Wirkung ist \emph{frei}: Fixiert $\pi$ jedes Tupelglied, so fixiert
$\pi$ jedes Bild punktweise, wegen der Ausschöpfung also ganz $[m]$,
d.\,h.\ $\pi=\mathrm{id}$. (Schablonen der Stelligkeit 0 stören
nicht: Sie tragen nichts zur Vereinigung bei.) Also $m!\mid N_m$,
und $(Q)^{\underline d}=\sum_m(N_m/m!)\,\ff Xm$ hat Koeffizienten in
$\Nn$ -- Positivität und Ganzzahligkeit in einem Zug.

Monotonie in $P$: $\Nn$-Linearkombination der
$(Q)^{\underline d}$. Monotonie in $Q$: Aus $a_k\le a_k'$ folgt
$T_k\subseteq T_k'$; jedes Tupel von $Q$-Strukturen ist eines von
$Q'$-Strukturen mit gleicher Bildvereinigung, also $N_m\le N_m'$
termweise. (Abschluss und Extremalprinzip rechnerisch exhaustiv
bestätigt für $r,s\le4$, der Abschluss zusätzlich für 300
Zufallspaare bis $\sigma=8$.)
\end{proof}

\begin{satz}[Shift-Überlauf; Register S42]
Für $E\,P(X)=P(X{+}1)$ gilt
$\Gamma_E(r)=r+\lfloor r^2/4\rfloor=\Gamma_\Delta(r{+}1)$.
\end{satz}

\begin{proof}
$E$ ist linear mit nichtnegativen Konstanten
($\ff{(X{+}1)}d=\ff Xd+d\,\ff X{d-1}$), also greift das
Extremalprinzip: $\Gamma_E(r)=\sigma(E\,M_r)$ mit Koeffizienten
$c_k=(r{-}k)+(k{+}1)(r{-}k{-}1)$ für $k\le r{-}2$ und $c_{r-1}=1$;
somit $k+c_k=r+(k{+}1)(r{-}k{-}1)$, und das Produkt zweier Zahlen
mit Summe $r$ ist maximal $\lfloor r^2/4\rfloor$. Die zweite
Gleichheit ist die Identität
$r+\lfloor r^2/4\rfloor=\lfloor(r{+}2)^2/4\rfloor-1$.
\end{proof}

\begin{satz}[Hori-\#P-Satz; Register S43]
Ist $f\in\#\mathrm P$ und $P$ eine Gestalt (nichtnegative
ganzzahlige Koeffizienten in der fallenden Fakultätsbasis), so ist
$P\circ f\in\#\mathrm P$.
\end{satz}

\begin{proof}
Sei $f(x)$ die Anzahl der Zeugen von $x$ unter einem polynomiellen
Verifizierer $V$, und $P=\sum_{d=0}^m c_d\ff Xd$. Neue Maschine auf
Eingabe $x$: Rate $d\in\{0,\ldots,m\}$, rate eine Variante
$\ell\in\{1,\ldots,c_d\}$, rate $d$ Zeugenkandidaten
$w_1,\ldots,w_d$; akzeptiere genau dann, wenn $V$ alle $w_i$
akzeptiert und die $w_i$ paarweise verschieden sind. Da $m$ und
alle $c_d$ feste Konstanten sind, bleibt die Maschine polynomiell,
und ihre Anzahl akzeptierender Pfade ist exakt
$\sum_d c_d\,f(x)^{\underline d}=P(f(x))$, denn
$f(x)^{\underline d}$ zählt die geordneten $d$-Tupel paarweise
verschiedener Zeugen. (Die Aussage folgt wegen
$\ff Xd=d!\binom Xd$ auch aus der bekannten Charakterisierung der
\emph{binomial-good}-Polynome als \#P-Abschlussoperationen,
Ikenmeyer--Pak, FOCS 2022; der obige Beweis ist selbständig.)
\end{proof}

\section*{Zu Kapitel 5: die beiden Halbringe und die Dualität}

\begin{satz}[Satz~\ref{satz:a3}, vollständig]
Mit $(n,x)\hplus(m,y)=(\max(n,m,\muh(x{+}y)),x{+}y)$ und
$(n,x)\htimes(m,y)=(\muh(xy),xy)$ ist $\Hori$ ein kommutativer
Halbring ohne Eins; $(0,0)$ ist neutral und absorbierend;
$h\htimes(1,1)=\KV(h)$; die wertkompakten Elemente bilden einen zu
$(\Nn,+,\cdot)$ isomorphen Unterhalbring mit Eins $(1,1)$.
\end{satz}

\begin{proof}
Vorbemerkung: $\muh$ ist monoton. \emph{Assoziativität von
$\hplus$:} Beide Klammerungen von $(n,x)$, $(m,y)$, $(k,z)$ liefern
den Wert $x{+}y{+}z$ und die Horizonte
$\max(n,m,k,\muh(x{+}y),\muh(x{+}y{+}z))$ bzw.\ mit $\muh(y{+}z)$;
wegen $x{+}y\le x{+}y{+}z$ und $y{+}z\le x{+}y{+}z$ absorbiert
$\muh(x{+}y{+}z)$ die inneren Terme. \emph{Assoziativität von
$\htimes$:} beide Seiten $(\muh(xyz),xyz)$.
\emph{Distributivität:} Links steht
$(n,x)\htimes\bigl((m,y)\hplus(k,z)\bigr)
=(\muh(xy{+}xz),\,xy{+}xz)$, rechts
$(\max(\muh(xy),\muh(xz),\muh(xy{+}xz)),\,xy{+}xz)$; wegen
$xy\le xy{+}xz$ und $xz\le xy{+}xz$ absorbiert $\muh(xy{+}xz)$
auch hier die inneren Terme.
\emph{Neutral/absorbierend:} $(n,x)\hplus(0,0)=(\max(n,0,\muh(x)),x)
=(n,x)$ wegen $n\ge\muh(x)$; $(n,x)\htimes(0,0)=(0,0)$.
\emph{Keine Eins:} Ein neutrales $(k,e)$ erforderte $xe=x$ für alle
$x$, also $e=1$, und $\muh(x)=n$ für alle Elemente -- das scheitert
an $(1,0)$. \emph{Projektor:}
$(n,x)\htimes(1,1)=(\muh(x),x)=\KV(n,x)$. \emph{Eck:} Für
wertkompakte Operanden ist
$\max(\muh(x),\muh(y),\muh(x{+}y))=\muh(x{+}y)$; unter $\htimes$ ist
das Ergebnis per Definition wertkompakt; $(1,1)$ ist dort neutral,
und $x\mapsto(\muh(x),x)$ ist der Isomorphismus.
\end{proof}

\begin{lemma}[$\sigma$-Additionsmonotonie]\label{lem:app-sigma}
$\sigma(P{+}Q)\ge\max(\sigma P,\sigma Q)$.
\end{lemma}

\begin{proof}
An der Maximalstelle $d^*$ von $P$ ist
$c_{d^*}(P{+}Q)\ge c_{d^*}(P)>0$, also
$\sigma(P{+}Q)\ge d^*+c_{d^*}(P{+}Q)\ge\sigma(P)$.
\end{proof}

\begin{satz}[Satz~\ref{satz:dualitaet}, vollständig]
Mit $(n,P)\boxplus(m,Q)=(\max(n,m,\sigmah(P{+}Q)),P{+}Q)$ und
$(n,P)\boxtimes(m,Q)=(\sigmah(PQ),PQ)$ ist $\Hori$ ein kommutativer
Halbring ohne Eins; $(0,0)$ ist neutral und absorbierend;
$h\boxtimes(1,1)=\KS(h)$; $\KS$ ist Homomorphismus beider
Operationen; die strukturkompakten Elemente bilden einen zu
$\Struk$ isomorphen Unterhalbring.
\end{satz}

\begin{proof}
Wörtlich spiegelbildlich zum vorigen Beweis, mit
Lemma~\ref{lem:app-sigma} anstelle der $\muh$-Monotonie. Zusätzlich
die Homomorphie von $\KS$: für $\boxplus$ ist
$\KS(a)\boxplus\KS(b)=(\max(\sigmah P,\sigmah Q,\sigmah(P{+}Q)),P{+}Q)
=(\sigmah(P{+}Q),P{+}Q)=\KS(a\boxplus b)$ nach dem Lemma; für
$\boxtimes$ sind beide Seiten $(\sigmah(PQ),PQ)$. Kein
$\boxtimes$-neutrales Element: es erforderte $PE=P$, also $E=1$,
und $\sigmah(P)=n$ für alle Elemente -- das scheitert an $(2,1)$,
denn $\sigmah(1)=1$.
\end{proof}

\begin{satz}[Produktprinzip, Satz~\ref{satz:produktprinzip}]
Jede kommutative Halbringstruktur $(\uplus,\odot)$ auf den
Exzessen liefert via
$(x,\varepsilon)\oplus(y,\delta)=(x{+}y,\varepsilon\uplus\delta)$,
$(x,\varepsilon)\otimes(y,\delta)=(xy,\varepsilon\odot\delta)$
eine gesetzestreue Arithmetik auf $\Hori$.
\end{satz}

\begin{proof}
Die Abbildung $(n,x)\mapsto(x,\,n-\muh(x))$ ist eine Bijektion
$\Hori\to\Nn\times\Nn$ (die Bedingung $n\ge\muh(x)$ ist genau
$\varepsilon\ge0$; die Umkehrung ist
$(x,\varepsilon)\mapsto(\muh(x){+}\varepsilon,x)$). Die Operationen
sind komponentenweise definiert, mit $(\Nn,+,\cdot)$ im Wertfaktor;
ein Produkt zweier kommutativer Halbringe erfüllt alle Gesetze
komponentenweise. Die Selektoren sind gültig, da der
Ergebnishorizont $\muh(\text{Wert})+\text{Exzess}\ge\muh$ ist.
\end{proof}

\begin{satz}[Kompaktifizierungsmonoid]
$\KS$ erhält Wertkompaktheit. Folglich gilt
$\KV\KS\KV=\KS\KV$, und das von $\KS,\KV$ erzeugte
Transformationsmonoid hat genau sechs Elemente -- $\mathrm{id}$,
$\KS$, $\KV$, $\KS\KV$, $\KV\KS$ und $\KS\KV\KS$ --, von denen
alle außer $\KV\KS$ idempotent sind.
\end{satz}

\begin{proof}
\emph{Stabilität:} Sei $(m,P)$ wertkompakt; das Element $(0,0)$
wird von $\KS$ fixiert, sei also $m\ge1$ und damit
$P(m{+}1)\ge m!$. Sei $s=\sigmah(P)$; angenommen, es gälte
$P(s{+}1)<s!$. Für $s\le t<m$ wächst
jeder Term $c_d\ff Xd$ von $t{+}1$ auf $t{+}2$ um den Faktor
$(t{+}2)/(t{+}2{-}d)$, wegen $d\le s{-}1\le t{-}1$ also um
höchstens $(t{+}2)/3$; damit
$P(t{+}2)<t!\,(t{+}2)/3\le(t{+}1)!$, und Induktion liefert
$P(m{+}1)<m!$ -- Widerspruch.
\emph{Relation:} Für jedes $h$ ist $\KV(h)$ wertkompakt, nach der
Stabilität auch $\KS(\KV(h))$, also fixiert $\KV$ dieses Element.
\emph{Sechs Elemente:} Jedes alternierende Wort der Länge $\ge4$
enthält $\KV\KS\KV$ und reduziert; Quadrate reduzieren per
Idempotenz -- höchstens sechs. Paarweise Verschiedenheit durch
explizite Trennzeugen (endliche Enumeration); insbesondere trennt
$h=(6,6)$: $\KV\KS(h)=(3,6)$, aber $\KS\KV\KS(h)=(2,5)$. Daraus
$(\KV\KS)^2=\KS\KV\KS\ne\KV\KS$: nicht idempotent; die Idempotenz
der übrigen fünf rechnet man mit der Relation nach.
\end{proof}

\section*{Zu Kapitel 5: das Tiefengesetz}

\begin{satz}[Satz~\ref{satz:tiefengesetz}, vollständig; Register S33]
Für jede Schale $m\ge3$ gilt
$\min\{\delta(h)\mid\muh(\val(h))=m\}=-g(m)$ mit
$g(m)=m-s_{\min}(m)$,
$s_{\min}(m)=\min\{s\mid M_s(m{+}1)\ge m!\}$; die
Terminaldarstellung von $x^*=M_{s_{\min}}(m{+}1)$ ist ein Zeuge.
\end{satz}

\begin{proof}
\emph{Schranke.} Sei $h=(n,x)$ mit $\muh(x)=m$, also $x\ge m!$.
Links: $\KV(h)=(m,x)$, und für $P=\operatorname{struct}(x,m)$ gilt
koeffizientenweise $P\preceq M_{\sigma(P)}$, also
$x=P(m{+}1)\le M_{\sigma(P)}(m{+}1)$; wegen $x\ge m!$ folgt
$\sigma(P)\ge s_{\min}$ und damit
$\hor(\KS\KV(h))=\sigma(P)\ge s_{\min}$. Rechts: Für
$P_n=\operatorname{struct}(x,n)$ ist die Auswertung monoton
(nichtnegative Koeffizienten), also
$\val(\KS(h))=P_n(\sigma(P_n){+}1)\le P_n(n{+}1)=x<(m{+}1)!$ und
damit $\hor(\KV\KS(h))\le m$. Zusammen
$\delta(h)\ge s_{\min}-m=-g(m)$.

\emph{Annahme.} Es ist $x^*=M_{s_{\min}}(m{+}1)\in[m!,(m{+}1)!)$
-- untere Grenze nach Definition von $s_{\min}$, obere wegen
$M_{s_{\min}}\preceq M_m$ und $M_m(m{+}1)=(m{+}1)!-1$. Sei
$h^*=(x^*,x^*)$ die Terminaldarstellung. $\KS$ fixiert $h^*$
(konstante Gestalt), also
$\hor(\KV\KS(h^*))=\muh(x^*)=m$. Andererseits ist $M_{s_{\min}}$
im Horizont $m$ eine gültige Ziffernfolge
($c_d\le s_{\min}-d\le m-d$) mit Wert $x^*$, nach der Bijektion
aus Kapitel~\ref{ch:horizonte} also \emph{die} Ziffernfolge:
$\operatorname{struct}(x^*,m)=M_{s_{\min}}$, und
$\hor(\KS\KV(h^*))=\sigma(M_{s_{\min}})=s_{\min}$. Also
$\delta(h^*)=s_{\min}-m=-g(m)$.
\end{proof}

\section*{Zu Kapitel 5: der Seitentreue-Satz}

\begin{satz}[Satz~\ref{satz:seitentreue}, Wert-Gestalt-Fälle]
\emph{(A)} Ist $\oplus$ wertgetragen und $\otimes$ strukturgetragen,
so ist bereits die Links-Distributivität für jede Wahl der
Horizontregeln unerfüllbar. \emph{(B)} Ist $\oplus$ strukturgetragen
und $\otimes$ wertgetragen, so sind Links-Distributivität und
Kommutativität von $\otimes$ zusammen unerfüllbar.
\end{satz}

\begin{proof}
\emph{(A).} Sei $u=(1,1)$ (Struktur $1$), $e_1=u$,
$e_m=e_{m-1}\oplus u$; da $\oplus$ wertgetragen ist, gilt
$\val(e_m)=m$, und die Struktur $Q_m$ von $e_m$ ist irgendeine
gültige Darstellung von $m$. Links-Distributivität liefert induktiv
\begin{equation}\label{eq:app-kette}
\val(a\otimes e_m)=m\cdot\val(a\otimes u).
\end{equation}
\emph{Schritt 1: $Q_m$ ist konstant.} Für $c\ge1$ hat $a_c=(c,c)$
die konstante Struktur $c$ (Terminaldarstellung); Konstanten werten
überall zu sich selbst aus, also $\val(a_c\otimes u)=c$. Mit
\eqref{eq:app-kette} folgt $c\,Q_m(y_c{+}1)=mc$, also
$Q_m(y_c{+}1)=m$, wobei $y_c$, der Horizont von $a_c\otimes e_m$,
wegen $y_c\ge\sigmah(cQ_m)=\max_d(d+c\,q_d)\ge c$ unbeschränkt
wächst ($Q_m\ne0$, denn $\val(e_m)=m\ge1$).
Alle Auswertungspunkte liegen im streng monotonen Bereich
(Argumente $\ge\sigmah+1\ge\deg+2$; fallende Fakultäten sind dort
streng wachsend). Bei $\deg Q_m\ge1$ nähme $Q_m$ den Wert $m$
höchstens einmal an -- Widerspruch. Also ist $Q_m$ konstant, und
mit $a=u$ folgt $Q_m=m$.
\emph{Schritt 2: Konstanten sterben.} Sei $g=(2,3)$ (Struktur $X$),
sei $t_g$ der Horizont von $g\otimes u$ und
$r=\val(g\otimes u)=t_g{+}1\ge\sigmah(X)+1=3$ fest. Nach
Schritt~1 hat $g\otimes e_m$ die Struktur $X\cdot m=mX$ mit
$\sigmah(mX)=m{+}1$; auf seinem Horizont $T_m\ge m{+}1$ ist also
$\val(g\otimes e_m)=m(T_m{+}1)\ge m(m{+}2)$; nach
\eqref{eq:app-kette} ist dieser Wert $mr$ -- Widerspruch für
$m>r-2$.

\emph{(B).} Nun addieren sich Strukturen und multiplizieren sich
Werte. Sei $f_1=u$, $f_m=f_{m-1}\oplus u$; dann hat $f_m$ die
konstante Struktur $m$ und den Wert $m$. Sei $A$ die Struktur von
$g\otimes u$. Induktiv hat $g\otimes f_m$ die Struktur $mA$; als
$\otimes$-Ergebnis vom Wert $3m$, realisiert auf seinem Horizont
$T_m$, gilt $\sigmah(mA)\le T_m$ nach der Gültigkeit der
Darstellung (Ziffernschranke), also $T_m\ge m$, und
$A(T_m{+}1)=3$ an unbeschränkt vielen Stellen. Wie in Schritt~1 ist
$A$ konstant, also $A=3$. Betrachte $f_2'=g\oplus g$ mit Struktur
$2X$, Horizont $H\ge\sigmah(2X)=3$, Wert $2(H{+}1)\ge8$.
Links-Distributivität und Kommutativität geben
$u\otimes(g\oplus g)=(g\otimes u)\oplus(g\otimes u)$ mit Struktur
$2\cdot3=6$ (Konstante), also Wert $6$; links steht aber
$1\cdot2(H{+}1)\ge8$. Widerspruch.
\end{proof}

\begin{satz}[nichtkommutativer Fall]
Ist $\oplus$ strukturgetragen und $\otimes$ wertgetragen und gelten
beide Distributivgesetze, so ist das System unerfüllbar -- ohne
Kommutativität, Assoziativität oder neutrale Elemente.
\end{satz}

\begin{proof}
Die Rechts-Distributivität liefert induktiv
$P_{f_m\otimes a}=m\cdot P_{u\otimes a}$; da $f_m\otimes a$ den
Wert $m\cdot\val(a)$ hat und der Horizont
$\ge\sigma(m\,P_{u\otimes a})\ge m$ wächst, ist $P_{u\otimes a}$
konstant gleich $\val(a)$ (Monotonie-Argument wie oben).
Insbesondere $P_{u\otimes g}=3$; das $g\oplus g$-Finale läuft dann
wörtlich wie im kommutativen Fall über die Links-Distributivität.
\end{proof}

\section*{Zu Kapitel 5 und 8: die dreiseitige Seitentreue}

\begin{satz}[Trilogie]
Bruchgetragene Addition ist nie total; wertgetragene Addition mit
bruchgetragener Multiplikation sowie strukturgetragene Addition mit
bruchgetragener Multiplikation sind für jede Selektorwahl
unerfüllbar (Links-Distributivität genügt).
\end{satz}

\begin{proof}
\emph{Totalität:} $F(u)+F(u)=\tfrac12+\tfrac12=1\notin[0,1)$, und
jedes Element hat Bruch $<1$.

\emph{Wert-$\oplus$, Bruch-$\otimes$:} Ketten $e_m$ mit
$\val(e_m)=m$, Horizont $N_m\ge\mu(m)$, Bruch $m/(N_m{+}1)!$.
Distributivität gibt $\val(u\otimes e_m)=m\,w$ mit
$w=\val(u\otimes u)$ -- ganzzahlig und $\ge1$, denn der Bruch von
$u\otimes u$ ist $\tfrac14>0$. Andererseits ist, mit $h_m$ dem
Horizont von $u\otimes e_m$,
$\val(u\otimes e_m)=\tfrac12\cdot\frac{m}{(N_m+1)!}(h_m{+}1)!$,
also $(h_m{+}1)!/(N_m{+}1)!=2w>1$ konstant -- aber ein
Fakultätsquotient $>1$ ist mindestens $N_m{+}2\to\infty$.

\emph{Struktur-$\oplus$, Bruch-$\otimes$:} Ketten $f_m$ mit
konstanter Struktur $m$, Horizont $N_m\ge m$, Bruch
$m/(N_m{+}1)!$. Distributivität gibt $P_{u\otimes f_m}=m\,A$ mit
$A=P_{u\otimes u}$, $\deg A=:d$; der Bruchvergleich auf dem
Horizont $T_m$ von $u\otimes f_m$ liefert nach Kürzen von $m$
\[
2\,A(T_m{+}1)=(T_m{+}1)^{\underline{\,k_m}},\qquad
k_m=T_m-N_m\ge1 .
\]
Das Produkt rechts ist mindestens $(N_m+k_m/2)^{k_m/2}$, die linke
Seite höchstens $C(N_m{+}k_m{+}1)^d$; wegen $N_m\to\infty$ ist
$k_m\le2d{+}1$ für große $m$, ein Wert $k^*$ tritt unendlich oft
auf, und $2A(Y)=Y^{\underline{k^*}}$ gilt als Polynomidentität.
Dann hätte $A$ den Koeffizienten $\tfrac12$ -- Ziffern sind ganz.
(Allgemeine Multiplikatoren $a$: dieselbe Rechnung mit
$c=(n_a{+}1)!/v_a>1$ statt $2$.)
\end{proof}

\section*{Zu Kapitel 6: Asymptotik und Ziffernformel}

\begin{satz}[Asymptotik der Sprungstellen]
Sei $m_g$ die kleinste Schale mit
$M_{m-g}(m{+}1)\ge m!$, wobei $M_s=\sum_{d<s}(s{-}d)\ff Xd$ das
Maximalpolynom ist. Mit
$c_g=(g{+}2)!\sum_{k\ge1}k/(g{+}k{+}1)!$ gilt
$(g{+}2)!/c_g\le m_g+1\le(g{+}2)!$ und $c_g=1+O(1/g)$, also
$m_g/(g{+}2)!\to1$.
\end{satz}

\begin{proof}
Umindexieren ($k=s{-}d$, $s=m{-}g$) liefert exakt
\[
\frac{M_{m-g}(m{+}1)}{m!}
=(m{+}1)\sum_{k=1}^{m-g}\frac{k}{(g{+}k{+}1)!}=:F(m).
\]
$F$ wächst in $m$ (der Vorfaktor streng, die Summe durch neu
hinzukommende nichtnegative Terme), die Schwelle ist also
wohldefiniert. Für $m{+}1\ge(g{+}2)!$ genügt der Term $k{=}1$:
$F(m)\ge(m{+}1)/(g{+}2)!\ge1$. Umgekehrt ist
$F(m)\le(m{+}1)\,c_g/(g{+}2)!$, also $F(m)<1$ für
$m{+}1<(g{+}2)!/c_g$.
Schließlich ist mit $x=1/(g{+}3)$
\[
c_g\le1+\sum_{k\ge2}k\,x^{k-1}
=1+\frac{x(2-x)}{(1-x)^2}=1+O(1/g). \qedhere
\]
\end{proof}

\begin{satz}[Ziffernformel für die Nicht-Springer]
Sei $d=(d_1,\ldots,d_n)$ die $H_n$-Ziffernfolge von $n!$ und $P$
die erste Position mit maximaler Ziffer ($d_P=P$), bzw.\ $P=n$,
falls keine existiert. Dann gilt
\[
N(n)=\sum_{p=1}^{P}\min(d_p,p)\cdot\frac{n!}{p!}.
\]
\end{satz}

\begin{proof}
$N(n)$ zählt Ziffernfolgen mit $a_i\le i-1$ und Wert $<n!$. Wir
zerlegen nach der ersten Position $p$, an der sich $a$ von $d$
unterscheidet. Der Präfix muss mit $d$ übereinstimmen -- möglich
nur, solange $d_j\le j-1$, also bis zur ersten maximalen Ziffer
einschließlich. An der Stelle $p$ muss $a_p<d_p$ und $a_p\le p-1$
gelten: $\min(d_p,p)$ Möglichkeiten. Der Suffix ist frei innerhalb
der Kappen: $\prod_{j>p}j=n!/p!$ Möglichkeiten, und jeder solche
Suffix bleibt strikt unterhalb, weil der volle (ungekappte)
Suffixraum das Intervall $[0,w_p)$ mit dem Stellengewicht
$w_p=(n{+}1)!/(p{+}1)!$ bijektiv füllt. Jeder Nicht-Springer wird genau
einmal gezählt.
\end{proof}

\begin{satz}[$\beta$-Reduktion]
$\ln\beta(r,r)=\ln\max_{t}\binom{r-1}{t}^2(r{-}1{-}t)!+O(\log r)$.
\end{satz}

\begin{proof}
Untere Schranke: Das Tripel $(p,q,j)=(r{-}1,r{-}1,r{-}1{-}t)$ trägt
zum Grad $r{-}1{+}t$ den Summanden
$\binom{r-1}{t}^2(r{-}1{-}t)!$ bei, und alle Beiträge sind
nichtnegativ. Obere Schranke: Jeder Beitrag ist höchstens
$r^2\binom{r-1}{j}^2j!$, und pro Grad gibt es höchstens $r^3$
Tripel; also $\beta\le2r+r^5\max_t(\cdot)$.
\end{proof}

\begin{satz}[$\beta$-Asymptotik]
$\ln(\beta(r,r)/r!)=2\sqrt r+O(\log r)$.
\end{satz}

\begin{proof}
Nach dem Reduktionssatz genügt die Aussage für den Maximalterm
$\binom st^2(s-t)!/r!=s!/(r\,t!^2(s-t)!)$ mit $s=r-1$. Nach oben:
Stirling-Schranken geben $\ln(s!/(t!^2(s-t)!))\le-h(t)+O(\log s)$
mit $h(t)=2t\ln t-t+(s{-}t)\ln(s{-}t)-s\ln s$; Konkavität liefert
$h(t)\ge t(2\ln t-\ln s-2)$, und diese Funktion hat ihr Minimum
$-2\sqrt s$ bei $t=\sqrt s$. Nach unten: $t=\lceil\sqrt s\rceil$
mit $\binom st\ge(s-t)^t/t!$ ergibt
$t\ln((s-t)/t^2)+2t-O(\log r)=2\sqrt s-O(\log r)$.
\end{proof}

\begin{satz}[Ein-Ziffern-Fakultäten]
$n!$ ist genau dann Einzelziffer in $H_n$, wenn $n=(i{+}1)!/d-1$
mit $1\le d\le i$; die Ziffer ist $d$ an Position $i$, und pro
Position gibt es genau $i$ Fälle.
\end{satz}

\begin{proof}
Einzelziffer $d$ an Position $i$ heißt $n!=d\,(n+1)!/(i+1)!$, also
$d=(i+1)!/(n+1)\le i$. Umgekehrt teilt jedes $d\le i$ die Zahl
$(i+1)!$, und $n+1=(i+1)!/d$ macht die Gleichung exakt; die
Position existiert, denn aus $n{+}1=(i{+}1)!/d\ge(i{+}1)!/i\ge
i{+}1$ folgt $i\le n$, und die eindeutige Ziffernentwicklung
besteht dann aus genau dieser Ziffer.
$d$ durchläuft bijektiv $1,\ldots,i$.
\end{proof}

\section*{Zu Kapitel 6 und 8: Interchange-Positivität, Scharfheit,
Nullzähler}

Für $x<(n{+}1)!$ sei $z_n(x):=P(n{+}2)$ mit
$P=\operatorname{struct}(x,n)$ -- der Wert nach einer
Nullerweiterung. Der \emph{Interchange-Defekt} eines Elements
$h=(n,x)$ ist
$\Omega(h)=z_n(x)-z_{n+1}(x)$: der Wert nach „erst wachsen, dann
werttreu liften“ minus der Wert nach „erst liften, dann wachsen“.

\begin{lemma}[Rekursionsformel]
Für $x=q(n{+}1)+r$, $0\le r\le n$:
$z_n(x)=r+(n{+}2)\,z_{n-1}(q)$.
\end{lemma}

\begin{proof}
Aus der Ziffernrekursion folgt
$P_{n,x}(X)=r+X\cdot P_{n-1,q}(X{-}1)$, denn der Koeffizientenshift
um einen Grad ist wegen
$\ff Xd=X\cdot\ff{(X{-}1)}{d-1}$ die Multiplikation mit $X$ bei
Argumentverschiebung. Auswerten bei $X=n{+}2$.
\end{proof}

\begin{lemma}[strikte Monotonie und Superadditivität]
$z_n(x)<z_n(x{+}1)$, folglich $z_n(y{+}s)\ge z_n(y)+s$.
\end{lemma}

\begin{proof}
Induktion über $n$; Basis $n=1$: $z_1(0)=0<1=z_1(1)$. Ohne
Übertrag wächst $z_n$ um genau 1; beim
Übertrag ($r=n$) ist
$z_n(x{+}1)-z_n(x)=(n{+}2)(z_{n-1}(q{+}1)-z_{n-1}(q))-n
\ge(n{+}2)-n=2$.
\end{proof}

\begin{lemma}[Multiplikationslemma]
Für $c\ge\ell{+}2$ und $cy<(\ell{+}1)!$ gilt
$z_\ell(cy)\ge(c{+}1)\,z_\ell(y)$.
\end{lemma}

\begin{proof}
Induktion über $\ell$; Basisfälle leer bzw.\ trivial. Sei
$y=q(\ell{+}1)+r$ und $cr=s(\ell{+}1)+u$ mit $0\le u\le\ell$; dann
$cy=(cq{+}s)(\ell{+}1)+u$ und nach der Rekursionsformel
$z_\ell(cy)=u+(\ell{+}2)\,z_{\ell-1}(cq{+}s)$. Superadditivität und
Induktionsvoraussetzung (Regime bleibt erhalten:
$c\ge\ell{+}2>(\ell{-}1){+}2$; Gültigkeit: $cq<\ell!$) geben
$z_{\ell-1}(cq{+}s)\ge(c{+}1)z_{\ell-1}(q)+s$. Es bleibt
$u+(\ell{+}2)s\ge(c{+}1)r$: Mit $c=\ell{+}1{+}e$, $e\ge1$, ist
$u=er\bmod(\ell{+}1)$, und die Ungleichung ist äquivalent zu
$r(c{-}\ell{-}1)=er\ge er\bmod(\ell{+}1)=u$ -- wahr.
\end{proof}

\begin{satz}[Interchange-Positivität]
$\Omega(h)\ge0$ für alle $h\in\Hori$.
\end{satz}

\begin{proof}
Nach Definition von $\Omega$ ist zu zeigen:
$z_{n+1}(x)\le z_n(x)$ für $x<(n{+}1)!$. Sei
$x=q'(n{+}2)+r'$. Dann
\[
z_{n+1}(x)=r'+(n{+}3)\,z_n(q')
\le r'+z_n\bigl((n{+}2)q'\bigr)
\le z_n\bigl(q'(n{+}2)+r'\bigr)=z_n(x),
\]
erst mit dem Multiplikationslemma ($c=n{+}2$, exakt an der
Regime-Grenze), dann mit der Superadditivität.
\end{proof}

\begin{satz}[Gleichheitsfälle; Register S32]
$\Omega(h)=0$ gilt genau dann, wenn beide Schritte dieser Kette
scharf sind: das Multiplikationslemma mit $c=n{+}2$ für den
\emph{scharfen Quotienten} $q'$, also
$z_n((n{+}2)q')=(n{+}3)\,z_n(q')$, und die Superadditivität
übertragsfrei, also $z_n((n{+}2)q'+r')=z_n((n{+}2)q')+r'$.
\end{satz}

\begin{proof}
$\Omega(h)$ ist die Summe der beiden Abweichungen
$z_n((n{+}2)q')-(n{+}3)z_n(q')$ und
$z_n((n{+}2)q'+r')-z_n((n{+}2)q')-r'$; beide sind nach
Multiplikationslemma bzw.\ Superadditivität nichtnegativ, ihre
Summe ist also genau dann null, wenn beide null sind.
\end{proof}

\begin{satz}[Ziffernentfaltung der Scharfheit]
Für $c\ge\ell+2$, $cy<(\ell+1)!$, $y=q(\ell+1)+r$, $e=c-\ell-1$
gilt: $z_\ell(cy)=(c{+}1)z_\ell(y)$ genau dann, wenn $er\le\ell$,
$(cq\bmod\ell)+r\le\ell-1$ und dieselbe Bedingung rekursiv für
$(\ell{-}1,c,q)$ gilt (Abbruch: für $\ell=0$ ist die Bedingung
leer).
\end{satz}

\begin{proof}
Die Kette $z_\ell(cy)=u+(\ell+2)z_{\ell-1}(cq+s)$ aus dem Beweis
des Multiplikationslemmas ist genau dann durchgehend scharf: die
arithmetische Abschätzung bei $er\le\ell$ (dann $s=r$, $u=er$),
die Superadditivität bei Übertragsfreiheit der $s=r$
Einzelschritte ab $cq$ -- deren Einerstelle ist $cq\bmod\ell$, die
Bedingung also $(cq\bmod\ell)+r\le\ell-1$ --, die dritte per
Definition rekursiv. Exhaustiv verifiziert für $\ell\le7$.
\end{proof}

\begin{satz}[Nullzähler-Schranke]
$\#\{\Omega=0\text{ auf }H_n\}\le(n{+}1)^{n+1}/n!$, der Anteil
fällt also mindestens wie $e^{-(1+o(1))n\ln n}$.
\end{satz}

\begin{proof}
Nach dem Gleichheitssatz ist jedes solche Element durch $(q',r')$
mit scharfem $q'$ bestimmt, und die Übertragsfreiheit erzwingt
$r'\le n$ (die Einerstelle wächst pro Schritt um eins und muss vor
jedem Schritt unter ihrem Maximum $n$ bleiben). Die Scharfheit
erzwingt nach der Ziffernentfaltung auf Ebene $\ell$ die
Ziffernschranke $r_\ell\le\lfloor\ell/(n{+}1{-}\ell)\rfloor$, und
$\lfloor\ell/j\rfloor+1\le(n{+}1)/j$ mit $j=n{+}1{-}\ell$ liefert
für die $q'$ das Produkt $(n{+}1)^n/n!$; zusammen mit den $n{+}1$
Möglichkeiten für $r'$ folgt die Schranke.
\end{proof}

\section*{Zu Kapitel 9: die Inversionsfolgen-Einordnung}

\begin{satz}[Register S44]
$H_n$ steht in Bijektion zu den Inversionsfolgen der Länge $n{+}1$
und zu $S_{n+1}$. Der Strukturhorizont ist
$\sigma(P_a)=n-\min_{e_j>0}((j{-}1)-e_j)$ für $a\ne0$ (für $a=0$
ist $\sigma=0$) mit $e=(0,a_1,\ldots,a_n)$, und
$|\{a\in H_n:\sigma(P_a)=r\}|=r\cdot r!$ für $1\le r\le n$.
\end{satz}

\begin{proof}
Setzt man $e_1=0$ und $e_{i+1}=a_i$, so ist $0\le e_j<j$ für alle
$j=1,\ldots,n{+}1$ -- die definierende Bedingung einer
Inversionsfolge; die Abbildung $\pi\mapsto(e_j)_j$ mit
$e_j=|\{i<j:\pi_i>\pi_j\}|$ ist die klassische Bijektion zu
$S_{n+1}$, und beide Seiten haben $(n{+}1)!$ Elemente. Für den
Strukturhorizont: $c_d=a_{n-d}=e_{n-d+1}$, also
$d+c_d=(n{+}1-j)+e_j$ mit $j=n-d+1$; Maximieren von $d+c_d$ über
$c_d>0$ ist Minimieren von $(j{-}1)-e_j$ über $e_j>0$, und
$\max(d+c_d)=n-\min_{e_j>0}((j{-}1)-e_j)$. Zählung: $\sigma(P_a)\le r$
bedeutet $a_i\le r-(n-i)$ für alle $i$ mit $a_i>0$, äquivalent
$a\le M$ koordinatenweise für das zu $H_r$ gehörige Maximalprofil;
die Zahl solcher $a$ in $H_n$ ist $(r{+}1)!$ (die Ziffern oberhalb
Position $r$ sind erzwungen null, darunter frei bis zur Schranke).
Also $|\{\sigma\le r\}|=(r{+}1)!$ und
$|\{\sigma=r\}|=(r{+}1)!-r!=r\cdot r!$. (Exhaustiv bestätigt bis
$n=6$, \texttt{experiments\_einordnung.py}.)
\end{proof}


\end{document}
